% !TEX root = ../Elements.tex
%%%%%%
% BOOK 9
%%%%%%
\pdfbookmark[0]{Book 9}{book9}
\pagestyle{empty}
\begin{center}
{\Huge ELEMENTS BOOK 9}\\
\spa\spa\spa
{\huge\it Applications of Number Theory\symbolfootnote[2]{The propositions contained in Books 7--9 are generally attributed to the
school of Pythagoras.}}
\end{center}
%\vspace{2cm}
%\epsfysize=3.5in
%\centerline{\epsffile{Book09/front.eps}}
\newpage

%%%%%%
% Prop 9.1
%%%%%%
\pdfbookmark[1]{Proposition 9.1}{pdf9.1}
\pagestyle{fancy}
\cfoot{{\greekfont τ̀ηεπαγε}}
\lhead{\large{\greekfont ΣΤΟΙΘΕΙΩΝ γ̀γν{9}.}}
\rhead{\large ELEMENTS BOOK 9}
\begin{Parallel}{}{} 
\ParallelLText{
\begin{center}
{\large \ggn{1}.}
\end{center}\vspace*{-7pt}

{\greekfont Ἐὰν δύο ὅμοιοι ἐπίπεδοι ἀριθμοὶ πολλαπλασιάσαντες ἀλλήλους
ποιῶσί τινα, ὁ γενόμενος τετράγωνος ἔσται.}\\

\epsfysize=1.2in
\centerline{\epsffile{Book09/fig01g.eps}}

{\greekfont Ἔστωσαν δύο ὅμοιοι ἐπίπεδοι ἀριθμοὶ οἱ Α, Β, καὶ ὁ Α τὸν
Β πολλαπλασιάσας τὸν Γ ποιείτω; λέγω, ὅτι ὁ Γ τετράγωνός ἐστιν.}

{\greekfont Ὁ γὰρ Α ἑαυτὸν πολλαπλασιάσας τὸν Δ ποιείτω. ὁ Δ ἄρα τετράγωνός ἐστιν. ἐπεὶ οὖν ὁ Α ἑαυτὸν μὲν πολλαπλασιάσας τὸν Δ πεποίηκεν, τὸν δὲ Β πολλαπλασιάσας τὸν Γ πεποίηκεν, ἔστιν ἄρα
ὡς ὁ Α πρὸς τὸν Β, οὕτως ὁ Δ πρὸς τὸν Γ. καὶ ἐπεὶ
οἱ Α, Β ὅμοιοι ἐπίπεδοί εἰσιν ἀριθμοί, τῶν Α, Β ἄρα εἷς
μέσος ἀνάλογον ἐμπίπτει ἀριθμός. ἐὰν δὲ δύο ἀριθμῶν
μεταχὺ κατὰ τὸ συνεθὲς ἀνάλογον ἐμπίπτωσιν ἀριθμοί, ὅσοι
εἰς α ὐτοὺς ἐμπίπτουσι, τοσοῦτοι καὶ εἰς τοὺς τὸν α ὐτὸν
λόγον ἔθοντας; ὥστε καὶ τῶν Δ, Γ εἷς μέσος ἀνάλογον
ἐμπίπτει ἀριθμός. καί ἐστι τετράγωνος ὁ Δ; τετράγωνος
ἄρα καὶ ὁ Γ; ὅπερ ἔδει δεῖχαι.}}

\ParallelRText{
\begin{center}
{\large Proposition 1}
\end{center}

If two similar plane numbers make some (number by) multiplying one another then the created (number) will be  square.\\

\epsfysize=1.2in
\centerline{\epsffile{Book09/fig01e.eps}}

Let $A$ and $B$ be two similar plane numbers, and let $A$ make $C$
(by) multiplying $B$. I say that $C$ is square.

For let $A$ make $D$ (by) multiplying itself. $D$ is thus square. Therefore, since $A$ has made $D$ (by) multiplying itself, and
has made $C$ (by) multiplying $B$, thus as $A$ is to $B$, so $D$ (is) to $C$ [Prop.~7.17]. And since $A$ and $B$ are similar
plane numbers, one number thus falls (between) $A$ and $B$ in mean
proportion [Prop.~8.18]. And if (some) numbers fall between two  numbers in continued proportion then, as many (numbers)
as fall in  (between) them (in continued proportion), so many also (fall) in (between numbers) having
the same ratio (as them in continued proportion)  [Prop.~8.8].
And hence one number falls (between) $D$ and $C$ in mean proportion.
And $D$ is  square. Thus, $C$ (is) also  square [Prop.~8.22]. (Which is)
the very thing it was required to show.}
\end{Parallel}

%%%%%%
% Prop 9.2
%%%%%%
\pdfbookmark[1]{Proposition 9.2}{pdf9.2}
\begin{Parallel}{}{} 
\ParallelLText{
\begin{center}
{\large \ggn{2}.}
\end{center}\vspace*{-7pt}

{\greekfont Ἐὰν δύο ἀριθμοὶ πολλαπλασιάσαντες ἀλλήλους ποιῶσι τετράγωνον,
ὅμοιοι ἐπίπεδοί εἰσιν ἀριθμοί.}

\epsfysize=1.2in
\centerline{\epsffile{Book09/fig02g.eps}}

{\greekfont Ἔστωσαν δύο ἀριθμοὶ οἱ Α, Β, καὶ ὁ Α τὸν Β πολλαπλασιάσας
τετράγωνον τὸν Γ ποιείτω; λέγω, ὅτι οἱ Α, Β ὅμοιοι
ἐπίπεδοί εἰσιν ἀριθμοί.}

{\greekfont Ὁ γὰρ Α ἑαυτὸν πολλαπλασιάσας τὸν Δ ποιείτω; ὁ Δ ἄρα
τετράγωνός ἐστιν. καὶ ἐπεὶ ὁ Α ἑαυτὸν μὲν πολλαπλασιάσας τὸν
Δ πεποίηκεν, τὸν δὲ Β πολλαπλασιάσας τὸν Γ πεποίηκεν, ἔστιν
ἄρα ὡς ὁ Α πρὸς  τὸν Β, ὁ Δ πρὸς τὸν Γ. καὶ ἐπεὶ ὁ Δ
τετράγωνός ἐστιν, ἀλλὰ καὶ ὁ Γ, οἱ Δ, Γ ἄρα ὅμοιοι
ἐπίπεδοί εἰσιν.
 τῶν Δ, Γ ἄρα εἷς μέσος ἀνάλογον
ἐμπίπτει. καί ἐστιν ὡς ὁ Δ πρὸς τὸν Γ, οὕτως ὁ Α πρὸς τὸν
Β; καὶ τῶν Α, Β ἄρα εῖ̔ς μέσος ἀνάλογον ἐμπίπτει.
ἐὰν δὲ δύο ἀριθμῶν εἷς μέσος ἀνάλογον ἐμπίπτῃ,
ὅμοιοι ἐπίπεδοί εἰσιν [οἱ] ἀριθμοί; οἱ ἄρα Α, Β ὅμοιοί
εἰσιν ἐπίπεδοι; ὅπερ ἔδει δεῖχαι.}}

\ParallelRText{
\begin{center}
{\large Proposition 2}
\end{center}

If two numbers make a square (number by)
multiplying one another then they are similar plane numbers.

\epsfysize=1.3in
\centerline{\epsffile{Book09/fig02e.eps}}

Let $A$ and $B$ be two numbers, and let $A$ make the square (number)
$C$ (by) multiplying $B$. I say that $A$ and $B$ are similar plane numbers.

For let $A$ make $D$ (by) multiplying itself. Thus, $D$ is square. And since $A$ has made $D$ (by) multiplying itself, and has made $C$ (by)
multiplying $B$, thus as $A$ is to $B$, so $D$ (is) to $C$ [Prop.~7.17]. And since $D$ is  square, and $C$ (is) also, $D$ and $C$ are thus similar plane numbers. Thus, 
one (number) falls (between) $D$ and $C$ in mean proportion [Prop.~8.18]. And as $D$ is to $C$, so $A$ (is) to $B$. Thus, one (number) also falls (between) $A$ and $B$ in mean proportion [Prop.~8.8].  And if one (number)
falls (between) two numbers in mean proportion then [the] numbers
are similar plane (numbers)  [Prop.~8.20].
Thus, $A$ and $B$ are similar plane (numbers). (Which is) the
very thing it was required to show.}
\end{Parallel}

%%%%%%
% Prop 9.3
%%%%%%
\pdfbookmark[1]{Proposition 9.3}{pdf9.3}
\begin{Parallel}{}{} 
\ParallelLText{
\begin{center}
{\large \ggn{3}.}
\end{center}\vspace*{-7pt}

{\greekfont Ἐὰν κύβος ἀριθμὸς ἑαυτὸν πολλαπλασιάσας ποιῇ τινα, ὁ γενόμενος
κύβος ἔσται.}

\epsfysize=1.2in
\centerline{\epsffile{Book09/fig03g.eps}}

{\greekfont Κύβος γὰρ ἀριθμὸς ὁ Α ἑαυτὸν πολλαπλασιάσας τὸν Β ποιείτω;
λέγω, ὅτι ὁ Β κύβος ἐστίν.}

{\greekfont Εἰλήφθω γὰρ τοῦ Α πλευρὰ ὁ Γ, καὶ ὁ Γ ἑαυτὸν πολλαπλασιάσας τὸν
Δ ποιείτω. φανερὸν δή ἐστιν, ὅτι ὁ Γ τὸν Δ πολλαπλασιάσας
τὸν Α πεποίηκεν. καὶ ἐπεὶ ὁ Γ ἑαυτὸν πολλαπλασιάσας τὸν Δ
πεποίηκεν, ὁ Γ ἄρα τὸν Δ μετρεῖ κατὰ τὰς ἐν α ὐτῷ μονάδας.
ἀλλὰ μὴν καὶ ἡ μονὰς τὸν Γ μετρεῖ κατὰ τὰς ἐν α ὐτῷ
μονάδας; ἔστιν ἄρα ὡς ἡ μονὰς πρὸς τὸν Γ, ὁ Γ πρὸς
τὸν Δ. πάλιν, ἐπεὶ ὁ Γ τὸν Δ πολλαπλασιάσας τὸν Α πεποίηκεν,
ὁ Δ ἄρα τὸν Α μετρεῖ κατὰ τὰς ἐν τῷ Γ μονάδας. μετρεῖ δὲ καὶ ἡ
μονὰς τὸν Γ κατὰ τὰς ἐν α ὐτῷ μονάδας; ἔστιν ἄρα ὡς ἡ μονὰς
πρὸς τὸν Γ, ὁ Δ πρὸς τὸν Α. ἀλλ ὡς ἡ μονὰς πρὸς τὸν Γ,
ὁ Γ πρὸς τὸν Δ; καὶ ὡς ἄρα ἡ μονὰς πρὸς τὸν Γ, οὕτως
ὁ Γ πρὸς τὸν Δ καὶ ὁ Δ πρὸς τὸν Α. τῆς ἄρα μονάδος
καὶ τοῦ Α ἀριθμοῦ δύο μέσοι ἀνάλογον
κατὰ τὸ συνεθὲς ἐμπεπτώκασιν ἀριθμοὶ οἱ Γ, Δ.
πάλιν, ἐπεὶ ὁ Α ἑαυτὸν πολλαπλασιάσας τὸν Β πεποίηκεν, ὁ Α
ἄρα τὸν Β μετρεῖ κατὰ τὰς ἐν α ὐτῷ μονάδας;
μετρεῖ δὲ καὶ ἡ μονὰς τὸν Α κατὰ τὰς ἐν α ὐτῷ
μονάδας;
 ἔστιν
ἄρα ὡς ἡ μονὰς πρὸς τὸν Α, ὁ Α πρὸς τὸν Β. τῆς δὲ μονάδος
καὶ τοῦ Α δύο μέσοι ἀνάλογον ἐμπεπτώκασιν ἀριθμοί; καὶ
τῶν Α, Β ἄρα δύο μέσοι ἀνάλογον ἐμπεσοῦνται ἀριθμοί. ἐὰν
δὲ δύο ἀριθμῶν δύο μέσοι ἀνάλογον ἐμπίπτωσιν, ὁ δὲ
πρῶτος κύβος ᾖ, καὶ ὁ δεύτερος κύβος ἔσται. καί ἐστιν
ὁ Α κύβος; καὶ ὁ Β ἄρα κύβος ἐστίν; ὅπερ ἔδει
δεῖχαι.}}

\ParallelRText{
\begin{center}
{\large Proposition 3}
\end{center}

If a cube number makes some (number by) multiplying itself then the created (number) will be cube.

\epsfysize=1.2in
\centerline{\epsffile{Book09/fig03e.eps}}

For let the cube number $A$ make $B$ (by) multiplying itself. I say
that $B$ is  cube.

For let the side $C$ of $A$ be taken. And let $C$ make
$D$ by multiplying itself. So it is clear that $C$ has made $A$ (by)
multiplying $D$. And since $C$ has made $D$ (by) multiplying itself,
$C$ thus measures $D$ according to the units in it [Def.~7.15]. But, in fact, a unit also measures
$C$ according to the units in it [Def.~7.20]. 
Thus, as a unit is to $C$, so $C$ (is) to $D$. Again, since $C$ has
made $A$ (by) multiplying $D$, $D$ thus measures $A$ according to
the units in $C$. And a unit also measures $C$ according to the units in it. Thus, as a unit is to $C$, so $D$ (is) to $A$. But, as a unit (is) to $C$, so
$C$ (is) to $D$. And thus as a unit (is) to $C$, so $C$ (is) to $D$, and
$D$ to $A$. Thus, two numbers, $C$ and $D$, have fallen (between) a unit
and the number $A$ in continued  mean proportion. Again, since $A$ has made $B$ (by) multiplying itself, $A$ thus measures $B$ according to the units in it.
And a unit also measures $A$ according to the units in it. Thus, as a
unit is to $A$, so $A$ (is) to $B$. And two numbers have
fallen (between) a unit and $A$ in mean proportion. Thus two
numbers will also fall (between) $A$ and $B$ in mean proportion [Prop.~8.8].
And if two
(numbers) fall (between) two numbers in mean proportion, and the
first (number) is cube, then the second will also be  cube 
[Prop.~8.23]. And $A$ is cube.
Thus, $B$ is also  cube. (Which is) the very thing it was required to show.}
\end{Parallel}

%%%%%%
% Prop 9.4
%%%%%%
\pdfbookmark[1]{Proposition 9.4}{pdf9.4}
\begin{Parallel}{}{} 
\ParallelLText{
\begin{center}
{\large \ggn{4}.}
\end{center}\vspace*{-7pt}

{\greekfont Ἐὰν κύβος ἀριθμὸς κύβον ἀριθμὸν πολλαπλασιάσας ποιῇ τινα,
ὁ γενόμενος κύβος ἔσται.}

{\greekfont Κύβος γὰρ ἀριθμὸς ὁ Α κύβον ἀριθμὸν τὸν Β πολλαπλασιάσας
τὸν Γ ποιείτω; λέγω, ὅτι ὁ Γ κύβος ἐστίν.}

\epsfysize=1.2in
\centerline{\epsffile{Book09/fig04g.eps}}

{\greekfont Ὁ γὰρ Α ἑαυτὸν πολλαπλασιάσας τὸν Δ ποιείτω;
ὁ Δ ἄρα κύβος ἐστίν. καὶ ἐπεὶ ὁ Α ἑαυτὸν μὲν
πολλαπλασιάσας τὸν Δ πεποίηκεν, τὸν δὲ Β πολλαπλασιάσας τὸν Γ
πεποίηκεν,  ἔστιν ἄρα ὡς ὁ Α πρὸς τὸν Β, οὕτως ὁ Δ πρὸς
τὸν Γ. καὶ ἐπεὶ οἱ Α, Β κύβοι εἰσίν, ὅμοιοι στερεοί εἰσιν οἱ Α, Β.
τῶν Α, Β ἄρα δύο μέσοι ἀνάλογον ἐμπίπτουσιν ἀριθμοί;
ὥστε καὶ τῶν Δ, Γ δύο μέσοι ἀνάλογον ἐμπεσοῦνται ἀριθμοί.
καί ἐστι κύβος ὁ Δ; κύβος ἄρα καὶ ὁ Γ; ὅπερ ἔδει δεῖχαι.}}

\ParallelRText{
\begin{center}
{\large Proposition 4}
\end{center}

If a cube number makes some (number by) multiplying a(nother) cube number then the created (number) will be  cube.

For let the cube number $A$ make $C$ (by) multiplying the cube number $B$. I say that $C$ is cube.

\epsfysize=1.2in
\centerline{\epsffile{Book09/fig04e.eps}}

For let $A$ make $D$ (by) multiplying itself.  Thus, $D$ is  cube 
[Prop.~9.3]. And since $A$ has made
$D$ (by) multiplying itself, and has made $C$ (by) multiplying $B$, 
thus as $A$ is to $B$, so $D$ (is) to $C$ [Prop.~7.17].  And since $A$ and $B$ are cube, $A$ and $B$ are similar solid (numbers). Thus, two numbers fall
(between) $A$ and $B$ in mean proportion [Prop.~8.19]. Hence, two numbers will
also fall (between) $D$ and $C$ in mean proportion [Prop.~8.8]. And $D$ is  cube. Thus,
$C$ (is) also cube  [Prop.~8.23]. 
(Which is) the very thing it was required to show.}
\end{Parallel}

%%%%%%
% Prop 9.5
%%%%%%
\pdfbookmark[1]{Proposition 9.5}{pdf9.5}
\begin{Parallel}{}{} 
\ParallelLText{
\begin{center}
{\large \ggn{5}.}
\end{center}\vspace*{-7pt}

{\greekfont Ἐὰν κύβος ἀριθμὸς ἀριθμόν τινα πολλαπλασιάσας κύβον ποιῇ,
καὶ ὁ πολλαπλασιασθεὶς κύβος ἔσται.}\\

\epsfysize=1.2in
\centerline{\epsffile{Book09/fig05g.eps}}

{\greekfont Κύβος γὰρ ἀριθμὸς ὁ Α ἀριθμόν τινα τὸν Β πολλαπλασιάσας κύβον
τὸν Γ ποιείτω; λέγω, ὅτι ὁ Β κύβος ἐστίν.}

{\greekfont Ὁ γὰρ Α ἑαυτὸν πολλαπλασιάσας τὸν Δ ποιείτω; κύβος ἄρα ἐστίν
ὁ Δ. καὶ ἐπεὶ ὁ Α ἑαυτὸν μὲν πολλαπλασιάσας τὸν Δ πεποίηκεν, τὸν
δὲ Β πολλαπλασιάσας τὸν Γ πεποίηκεν, ἔστιν ἀρα ὡς ὁ Α πρὸς τὸν Β,
ὁ Δ πρὸς τὸν Γ. καὶ ἐπεὶ οἱ Δ, Γ κύβοι εἰσίν, ὅμοιοι στερεοί
εἰσιν. τῶν Δ, Γ ἄρα δύο μέσοι ἀνάλογον ἐμπίπτουσιν ἀριθμοί.
καί ἐστιν ὡς ὁ Δ πρὸς τὸν Γ, οὕτως ὁ Α πρὸς τὸν Β; καὶ τῶν
Α, Β ἄρα δύο μέσοι ἀνάλογον ἐμπίπτουσιν ἀριθμοί. καί ἐστι
κύβος ὁ Α; κύβος ἄρα ἐστὶ καὶ ὁ Β; ὅπερ ἔδει δεῖχαι.}}

\ParallelRText{
\begin{center}
{\large Proposition 5}
\end{center}

If a cube number makes a(nother) cube number
(by) multiplying some (number) then the (number) multiplied will also be 
cube.

\epsfysize=1.2in
\centerline{\epsffile{Book09/fig05e.eps}}

For let the cube number $A$ make the cube (number) $C$ (by) multiplying
some number $B$. I say that $B$ is cube.

For let $A$ make $D$ (by) multiplying itself. $D$ is thus  cube [Prop.~9.3]. And since $A$ has made $D$ (by)
multiplying itself, and has made $C$ (by) multiplying $B$, thus as $A$ is to
$B$, so $D$ (is) to $C$ [Prop.~7.17]. And since
$D$ and $C$ are (both) cube, they are similar solid (numbers). Thus,
two numbers fall (between) $D$ and $C$ in mean proportion [Prop.~8.19]. And as $D$  is to $C$, so $A$ (is) to $B$.  Thus, two numbers also fall (between) $A$ and $B$ in mean proportion
[Prop.~8.8]. And $A$ is  cube. Thus,
 $B$ is also  cube  [Prop.~8.23]. (Which is) the very thing it was required to show.}
\end{Parallel}

%%%%%%
% Prop 9.6
%%%%%%
\pdfbookmark[1]{Proposition 9.6}{pdf9.6}
\begin{Parallel}{}{} 
\ParallelLText{
\begin{center}
{\large \ggn{6}.}
\end{center}\vspace*{-7pt}

{\greekfont Ἐὰν ἀριθμὸς ἑαυτὸν πολλαπλασιάσας κύβον ποιῇ, καὶ α ὐτὸς κύβος
ἔσται.}

{\greekfont Ἀριθμὸς γὰρ ὁ Α ἑαυτὸν πολλαπλασιάσας κύβον τὸν Β ποιείτω;
λέγω, ὅτι καὶ ὁ Α κύβος ἐστίν.}

\epsfysize=0.8in
\centerline{\epsffile{Book09/fig06g.eps}}

{\greekfont Ὁ γὰρ Α τὸν Β πολλαπλασιάσας τὸν Γ ποιείτω. ἐπεὶ οὖν ὁ Α
ἑαυτὸν μὲν πολλαπλασιάσας τὸν Β πεποίηκεν, τὸν δὲ Β πολλαπλασιάσας
τὸν Γ πεποίηκεν, ὁ Γ ἄρα κύβος ἐστίν. καὶ ἐπεὶ ὁ Α ἑαυτὸν
πολλαπλασιάσας τὸν Β πεποίηκεν, ὁ Α ἄρα τὸν Β μετρεῖ κατὰ τὰς
ἐν α ὐτῷ μονάδας. μετρεῖ δὲ καὶ ἡ μονὰς τὸν Α κατὰ τὰς
ἐν α ὐτῷ μονάδας. ἔστιν ἄρα ὡς ἡ μονὰς πρὸς τὸν Α, οὕτως
ὁ Α πρὸς τὸν Β. καὶ ἐπεὶ ὁ Α τὸν Β πολλαπλασιάσας τὸν Γ πεποίηκεν,
ὁ Β ἄρα τὸν Γ μετρεῖ κατὰ τὰς ἐν τῷ Α μονάδας. μετρεῖ δὲ
καὶ ἡ μονὰς  τὸν Α κατὰ τὰς ἐν α ὐτῷ μονάδας. ἔστιν ἄρα  ὡς ἡ μονὰς πρὸς τὸν Α,
οὕτως ὁ Β πρὸς
τὸν Γ. ἀλλ ὡς ἡ μονὰς πρὸς τὸν Α, οὕτως ὁ Α πρὸς τὸν Β;
καὶ ὡς ἄρα ὁ Α πρὸς τὸν Β, ὁ Β πρὸς τὸν Γ. καὶ ἐπεὶ οἱ Β, Γ
κύβοι εἰσίν, ὅμοιοι στερεοί εἰσιν. τῶν Β, Γ ἄρα δύο μέσοι ἀνάλογόν εἰσιν ἀριθμοί. καί ἐστιν ὡς ὁ Β πρὸς τὸν Γ, ὁ Α πρὸς
τὸν Β. καὶ τῶν Α, Β ἄρα δύο μέσοι ἀνάλογόν εἰσιν ἀριθμοί.
καί ἐστιν κύβος ὁ Β; κύβος ἄρα ἐστὶ καὶ ὁ Α; ὅπερ ἔδει
δεῖχαι.}}

\ParallelRText{
\begin{center}
{\large Proposition 6}
\end{center}

If a number makes a cube (number by) multiplying itself then it itself will also be  cube.

For let the number $A$ make the cube (number) $B$ (by) multiplying itself.
I say that $A$ is also cube.

\epsfysize=0.8in
\centerline{\epsffile{Book09/fig06e.eps}}

For let $A$ make $C$ (by) multiplying $B$. Therefore, since $A$
has made $B$ (by) multiplying itself, and has made $C$ (by) multiplying 
$B$, $C$ is thus cube. And since $A$ has made $B$ (by)
multiplying itself, $A$ thus measures $B$ according to the units in ($A$).
And a unit also measures $A$ according to the units in it. Thus, as
a unit is to $A$, so $A$ (is) to $B$. And since $A$ has made $C$ (by)
multiplying $B$, $B$ thus measures $C$ according to the units in $A$.
And a unit also measures $A$ according to the units in it. Thus, as
a unit is to $A$, so $B$ (is) to $C$. But, as a unit (is) to $A$, so $A$ (is)
to $B$. And thus as $A$ (is) to $B$, (so) $B$ (is) to $C$. And since $B$
and $C$ are cube, they are similar solid (numbers). Thus,
there exist two numbers in mean proportion (between) $B$ and $C$ [Prop.~8.19]. And as $B$ is to $C$, (so) $A$ (is)
to $B$. Thus, there also exist two numbers in mean proportion (between) $A$ and
$B$ [Prop.~8.8].  And $B$ is  cube. Thus, $A$
is also  cube   [Prop.~8.23]. (Which is)
the very thing it was required to show.}
\end{Parallel}

%%%%%%
% Prop 9.7
%%%%%%
\pdfbookmark[1]{Proposition 9.7}{pdf9.7}
\begin{Parallel}{}{} 
\ParallelLText{
\begin{center}
{\large \ggn{7}.}
\end{center}\vspace*{-7pt}

{\greekfont Ἐὰν σύνθετος ἀριθμὸς ἀριθμόν τινα πολλαπλασιά-σας ποιῇ τινα, ὁ
γενόμενος στερεὸς ἔσται.}\\

\epsfysize=1.5in
\centerline{\epsffile{Book09/fig07g.eps}}

{\greekfont Σύνθετος γὰρ ἀριθμὸς ὁ Α ἀριθμόν τινα τὸν Β πολλαπλασιάσας 
τὸν Γ ποιείτω; λέγω, ὅτι ὁ Γ στερεός ἐστιν.}

{\greekfont Ἐπεὶ γὰρ ὁ Α σύνθετός ἐστιν, ὑπὸ ἀριθμοῦ τινος μετρηθήσεται.
μετρείσθω ὑπὸ τοῦ Δ, καὶ ὁσάκις ὁ Δ τὸν Α μετρεῖ, τοσαῦται
μονάδες ἔστωσαν ἐν τῷ Ε. ἐπεὶ οὖν ὁ Δ τὸν Α μετρεῖ
κατὰ τὰς ἐν τῷ Ε μονάδας, ὁ Ε ἄρα τὸν Δ πολλαπλασιάσας τὸν Α
πεποίηκεν. καὶ ἐπεὶ ὁ Α τὸν Β πολλαπλασιάσας τὸν Γ πεποίηκεν,
ὁ δὲ Α ἐστιν ὁ ἐκ τῶν Δ, Ε, ὁ ἄρα ἐκ τῶν Δ, Ε τὸν Β
πολλαπλασιάσας τὸν Γ πεποίηκεν. ὁ Γ ἄρα στερεός ἐστιν, πλευραὶ
δὲ α ὐτοῦ εἰσιν οἱ Δ, Ε, Β; ὅπερ ἔδει δεῖχαι.}}

\ParallelRText{
\begin{center}
{\large Proposition 7}
\end{center}

If a composite number makes some (number by)
multiplying some (other) number then the created (number) will be  solid.

\epsfysize=1.5in
\centerline{\epsffile{Book09/fig07e.eps}}

For let the composite number $A$ make $C$ (by) multiplying some number
$B$. I say that $C$ is  solid.

For since $A$ is a composite (number), it will be measured by some
number. Let it be measured by $D$.  And, as many times as $D$ measures
$A$, so many units let there be in $E$. Therefore, since $D$ measures $A$
according to the units in $E$, $E$ has thus made $A$ (by) multiplying
$D$ [Def.~7.15]. And since $A$ has made $C$ (by) multiplying $B$, and $A$ is
the (number created) from (multiplying) $D$, $E$, the (number created)
from (multiplying) $D$, $E$ has thus made $C$ (by) multiplying $B$. Thus,
$C$ is solid, and its sides are $D$, $E$, $B$. (Which is) the
very thing it was required to show.}
\end{Parallel}

%%%%%%
% Prop 9.8
%%%%%%
\pdfbookmark[1]{Proposition 9.8}{pdf9.8}
\begin{Parallel}{}{} 
\ParallelLText{
\begin{center}
{\large \ggn{8}.}
\end{center}\vspace*{-7pt}

{\greekfont Ἐὰν ἀπὸ μονάδος ὁποσοιοῦν ἀριθμοὶ ἑχῆς ἀνάλογον ὦσιν,
ὁ μὲν τρίτος ἀπὸ τῆς μονάδος τετράγωνος ἔσται καὶ οἱ ἕνα
διαλείποντες, ὁ δὲ τέταρτος κύβος καὶ οἱ δύο διαλείποντες πάντες,
ὁ δὲ ἕβδομος κύβος ἅμα καὶ τετράγωνος καὶ οἱ πέντε
διαλείποντες.}\\~\\~\\

\epsfysize=1.8in
\centerline{\epsffile{Book09/fig08g.eps}}

{\greekfont Ἔστωσαν ἀπὸ μονάδος ὁποσοιοῦν ἀριθμοὶ ἑχῆς ἀνάλογ-ον
οἱ Α, Β, Γ, Δ, Ε, Ζ; λέγω, ὅτι ὁ μὲν τρίτος ἀπὸ τῆς
μονάδος ὁ Β τετράγωνός ἐστι καὶ οἱ ἕνα διαλείποντες πάντες,
ὁ δὲ τέταρτος ὁ Γ κύβος καὶ οἱ δύο διαλείποντες πάντες,
ὁ δὲ
ἕβδομος ὁ Ζ κύβος ἅμα καὶ τετράγωνος καὶ οἱ πέντε
διαλείποντες πάντες.}

{\greekfont Ἐπεὶ γάρ ἐστιν ὡς ἡ μονὰς πρὸς τὸν Α, οὕτως ὁ Α πρὸς τὸν Β,
ἰσάκις ἄρα ἡ μονὰς τὸν Α ἀριθμὸν μετρεῖ καὶ ὁ Α τὸν Β. ἡ δὲ
μονὰς τὸν Α ἀριθμὸν μετρεῖ κατὰ τὰς ἐν α ὐτῷ μονάδας; καὶ ὁ Α
ἄρα τὸν Β μετρεῖ κατὰ τὰς ἐν τῷ Α μονάδας. ὁ Α ἄρα ἑαυτὸν
πολλαπλασιάσας τὸν Β πεποίηκεν; τετράγωνος ἄρα ἐστὶν ὁ Β. καὶ ἐπεὶ
οἱ Β, Γ, Δ ἑχῆς ἀνάλογόν εἰσιν, ὁ δὲ Β τετράγωνός ἐστιν,
καὶ ὁ Δ ἄρα τετράγωνός ἐστινμήιὰ τὰ α ὐτὰ δὴ καὶ ὁ Ζ
τετράγωνός ἐστιν. ὁμοίως δὴ δείχομεν, ὅτι καὶ οἱ ἕνα
διαλείποντες πάντες τετράγωνοί εἰσιν. λέγω δή, ὅτι καὶ ὁ τέταρτος
ἀπὸ τῆς μονάδος ὁ Γ κύβος ἐστὶ καὶ οἱ δύο διαλείποντες πάντες.
ἐπεὶ γάρ ἐστιν ὡς ἡ μονὰς πρὸς τὸν Α, οὕτως ὁ Β πρὸς τὸν
Γ, ἰσάκις ἄρα ἡ μονὰς τὸν Α ἀριθμὸν μετρεῖ καὶ ὁ Β τὸν Γ.
ἡ δὲ μονὰς τὸν Α ἀριθμὸν μετρεῖ κατὰ τὰς ἐν τῷ Α μονάδας;
καὶ ὁ Β  ἄρα τὸν Γ μετρεῖ κατὰ τὰς ἐν τῷ Α μονάδας; ὁ Α
ἄρα τὸν Β πολλαπλασιάσας τὸν Γ πεποίηκεν. ἐπεὶ οὖν ὁ Α
ἑαυτὸν μὲν πολλαπλασιάσας τὸν Β πεποίηκεν, τὸν δὲ Β πολλαπλασιάσας
τὸν Γ πεποίηκεν, κύβος ἄρα ἐστὶν ὁ Γ. καὶ ἐπεὶ οἱ Γ, Δ, Ε, Ζ
ἑχῆς ἀνάλογόν εἰσιν, ὁ δὲ Γ κύβος ἐστίν, καὶ ὁ Ζ ἄρα κύβος
ἐστίν. ἐδείθθη δὲ καὶ τετράγωνος; ὁ ἄρα ἕβδομος ἀπὸ
τῆς μονάδος κύβος τέ ἐστι καὶ τετράγωνος. ὁμοίως δὴ δείχομεν,
ὅτι καὶ οἱ πέντε διαλείποντες πάντες κύβοι τέ εἰσι καὶ
τετράγωνοι; ὅπερ ἔδει δεῖχαι.}}

\ParallelRText{
\begin{center}
{\large Proposition 8}
\end{center}

If any multitude whatsoever of numbers is continuously proportional, (starting) from a unit, then the third from the unit will be 
square, and (all) those (numbers after that) which leave an interval of  one (number), and the fourth (will
be) cube, and all those (numbers after that) which leave an interval of two (numbers), and the seventh
(will be) both cube and square, and (all) those (numbers after that) which leave an interval of five (numbers).

\epsfysize=1.8in
\centerline{\epsffile{Book09/fig08e.eps}}

Let any multitude whatsoever of numbers, $A$, $B$, $C$, $D$, $E$, $F$, be continuously proportional, (starting) from a unit. I say that the third from the unit, $B$, is square, and all those (numbers after that) which leave an interval of one (number). And the fourth (from the unit), $C$,
(is)  cube, and all those (numbers after that) which leave an
interval of two (numbers). And the seventh (from the unit), $F$, (is)
both cube and square,  and all those (numbers after that) which leave an interval of five (numbers).

For since as the unit is to $A$, so $A$ (is) to $B$, the unit thus measures the number $A$
the same number of times as $A$ (measures) $B$ [Def.~7.20]. And the unit measures the number $A$
according to the units in it. Thus, $A$ also measures $B$ according to the units
in $A$. $A$ has thus made $B$ (by) multiplying itself [Def.~7.15].  Thus, $B$ is  square. And
since $B$, $C$, $D$ are continuously proportional, and $B$ is  square, $D$ is thus also  square  
[Prop.~8.22]. So, for the same (reasons), $F$ is also
 square. So, similarly, we can also show that all those (numbers after that) which leave an interval of  one (number) are square. 
So I also say that the fourth (number) from the unit, $C$,  is  cube, and
all those (numbers after that) which leave an interval of two (numbers).
For since as the unit is to $A$, so $B$ (is) to $C$, the unit thus measures the
number $A$ the same number of times that $B$ (measures) $C$. And the
unit measures the number $A$ according to the units in $A$. And thus
$B$ measures $C$ according to the units in $A$.  $A$ has thus made $C$
(by) multiplying $B$. Therefore, since $A$ has made $B$ (by) multiplying
itself, and has made $C$ (by) multiplying $B$, $C$ is thus  cube. And since $C$, $D$, $E$, $F$ are continuously proportional, and
$C$ is cube, $F$ is thus also  cube  [Prop.~8.23]. And it was also shown (to be) 
square. Thus, the seventh (number) from the unit is (both) 
cube and  square. So, similarly, we can show that  all those (numbers after that) which leave an interval of five (numbers) are (both)
cube and square. (Which is) the very thing it was required to show.}
\end{Parallel}

%%%%%%
% Prop 9.9
%%%%%%
\pdfbookmark[1]{Proposition 9.9}{pdf9.9}
\begin{Parallel}{}{} 
\ParallelLText{
\begin{center}
{\large \ggn{9}.}
\end{center}\vspace*{-7pt}

{\greekfont Ἐὰν ἀπὸ μονάδος ὁποσοιοῦν ἑχῆς κατὰ τὸ συνεθὲς ἀριθμοὶ
ἀνάλογον ὦσιν, ὁ δὲ μετὰ τὴν μονάδα τετράγωνος ᾖ, καὶ οἱ
λοιποὶ πάντες τετράγωνοι ἔσονται. καὶ ἐὰν ὁ μετὰ τὴν μονάδα
κύβος ᾖ, καὶ οἱ λοιποὶ πάντες κύβοι ἔσονται.}\\

\epsfysize=1.8in
\centerline{\epsffile{Book09/fig08g.eps}}

{\greekfont Ἔστωσαν ἀπὸ μονάδος ἑχῆς ἀνάλογον ὁσοιδηποτοῦν ἀριθμοὶ
οἱ Α, Β, Γ, Δ, Ε, Ζ,  ὁ δὲ μετὰ τὴν μονάδα ὁ Α τετράγωνος
ἔστω; λέγω, ὅτι καὶ οἱ λοιποὶ πάντες τετράγωνοι ἔσονται.}

{\greekfont Ὅτι μὲν οὖν ὁ τρίτος ἀπὸ τῆς μονάδος ὁ Β τετράγωνός
ἐστι καὶ οἱ ἕνα διαπλείποντες πάντες, δέδεικται; λέγω [δή], ὅτι
καὶ οἱ λοιποὶ πάντες τετράγωνοί εἰσιν. ἐπεὶ γὰρ οἱ Α, Β, Γ
ἑχῆς ἀνάλογόν εἰσιν, καί ἐστιν ὁ Α τετράγωνος, καὶ ὁ Γ [ἄρα]
τετράγωνος ἐστιν. πάλιν, ἐπεὶ [καὶ] οἱ Β, Γ, Δ ἑχῆς ἀνάλογόν
εἰσιν, καί ἐστιν ὁ Β τετράγωνος, καὶ ὁ Δ [ἄρα] τετράγωνός
ἐστιν. ὁμοίως δὴ δείχομεν, ὅτι καὶ οἱ λοιποὶ πάντες
τετράγωνοί εἰσιν.}

{\greekfont Ἀλλὰ δὴ ἔστω ὁ Α κύβος; λέγω, ὅτι καὶ οἱ λοιποὶ πάντες
κύβοι εἰσίν.}

{\greekfont Ὅτι μὲν οὖν ὁ τέταρτος ἀπὸ τῆς μονάδος ὁ Γ κύβος
ἐστὶ καὶ οἱ δύο διαλείποντες πάντες, δέδεικται;
λέγω [δή], ὅτι καὶ οἱ λοιποὶ πάντες κύβοι εἰσίν. ἐπεὶ
γάρ ἐστιν ὡς ἡ μονὰς πρὸς τὸν Α, οὕτως ὁ Α πρὸς
τὸν Β, ἰσάκις ἀρα ἡ μονὰς τὸν Α μετρεῖ καὶ ὁ Α τὸν Β.
ἡ δὲ μονὰς τὸν Α μετρεῖ κατὰ τὰς ἐν α ὐτῷ
μονάδας; καὶ ὁ Α ἄρα τὸν Β μετρεῖ κατὰ τὰς ἐν α ὐτῷ
μονάδας; ὁ Α ἄρα ἑαυτὸν πολλαπλασιάσας τὸν Β πεποίηκεν. καί
ἐστιν ὁ Α κύβος. ἐὰν δὲ κύβος ἀριθμὸς ἑαυτὸν πολλαπλασιάσας
ποιῇ τινα, ὁ γενόμενος κύβος ἐστίν; καὶ ὁ Β ἄρα κύβος ἐστίν.
καὶ ἐπεὶ τέσσαρες ἀριθμοὶ οἱ Α, Β, Γ, Δ ἑχῆς ἀνάλογόν
εἰσιν, καί ἐστιν ὁ Α κύβος, καὶ ὁ Δ ἄρα κύβος ἐστίν. διὰ
τὰ α ὐτὰ δὴ καὶ ὁ Ε κύβος ἐστίν, καὶ ὁμοίως οἱ λοιποὶ
πάντες κύβοι εἰσίν; ὅπερ ἔδει δεῖχαι.}}

\ParallelRText{
\begin{center}
{\large Proposition 9}
\end{center}

If  any multitude whatsoever of numbers is continuously proportional, (starting) from a unit, and the (number) after the
unit is square, then all the remaining (numbers) will also 
be square. And if the (number) after the unit is  cube, then
all the remaining (numbers) will also be cube.

\epsfysize=1.8in
\centerline{\epsffile{Book09/fig08e.eps}}

Let any multitude whatsoever of numbers, $A$, $B$, $C$, $D$, $E$, $F$, be continuously proportional, (starting) from a unit. And let the (number) after the
unit, $A$, be square. I say that all the remaining (numbers)
will also be square.

In fact, it has (already) been shown that the third (number) from the unit, $B$, is
square, and all those (numbers after that) which leave an interval of one (number) [Prop.~9.8]. [So] I say that all the remaining (numbers) are also square. For since $A$, $B$, $C$
are continuously proportional, and $A$ (is)  square, $C$
is [thus] also  square [Prop.~8.22]. Again, 
since $B$, $C$, $D$ are [also] continuously proportional, and $B$ is
 square, $D$ is [thus] also square [Prop.~8.22]. So, similarly,
we can show that all the remaining (numbers) are also square.

And so let $A$ be  cube. I say that all the remaining (numbers)
are also cube.

In fact, it has (already) been shown that the fourth (number) from the unit, $C$, is
cube, and all those (numbers after that) which leave an interval of two (numbers) [Prop.~9.8]. [So] I say
that all the remaining (numbers) are also cube. For since as the
unit is to $A$, so $A$ (is) to $B$, the unit thus measures $A$ the same number of times as $A$ (measures) $B$. And the unit measures $A$ according to the units in it. Thus, $A$ also measures $B$ according to the
units in ($A$).
$A$ has thus made $B$ (by) multiplying itself. 
And $A$ is  cube. And if a cube number makes some (number by)
multiplying itself then the created (number) is cube  [Prop.~9.3]. Thus, $B$ is also cube.
And since the four numbers $A$, $B$, $C$, $D$ are continuously proportional, and $A$ is  cube, $D$ is thus also  cube 
 [Prop.~8.23]. So, for the same (reasons), $E$ is also
 cube, and, similarly, all the remaining (numbers) are cube. (Which is) the very thing it was required to show.}
\end{Parallel}

%%%%%%
% Prop 9.10
%%%%%%
\pdfbookmark[1]{Proposition 9.10}{pdf9.10}
\begin{Parallel}{}{} 
\ParallelLText{
\begin{center}
{\large \ggn{10}.}
\end{center}\vspace*{-7pt}

{\greekfont Ἐὰν ἀπὸ μονάδος ὁποσοιοῦν ἀριθμοὶ [ἑχῆς] ἀνάλογ-ον ὦσιν,
ὁ δὲ μετὰ τὴν μονάδα μὴ ᾖ τετράγωνος, οὐδ ἄλλος οὐδεὶς
τετράγωνος ἔσται θωρὶς τοῦ τρίτου ἀπὸ τῆς μονάδος
καὶ τῶν ἕνα διαλειπόντων πάντων. καὶ ἐὰν ὁ μετὰ τὴν μονάδα κύβος μὴ ᾖ, οὐδὲ ἄλλος οὐδεὶς κύβος ἔσται θωρὶς τοῦ
τετάρτου ἀπὸ τῆς μονάδος καὶ τῶν δύο διαλειπόντων πάντων.}\\~\\~\\

\epsfysize=1.8in
\centerline{\epsffile{Book09/fig08g.eps}}

{\greekfont Ἔστωσαν ἀπὸ μονάδος ἑχῆς ἀνάλογον ὁσοιδηποτοῦν ἀριθμοὶ
οἱ Α, Β, Γ, Δ, Ε, Ζ, ὁ μετὰ τὴν μονάδα ὁ Α μὴ ἔστω τετράγωνος;
λέγω, ὅτι οὐδὲ ἄλλος οὐδεὶς τετράγωνος ἔσται θωρὶς τοῦ
τρίτου ἀπὸ τὴς μονάδος [καὶ τῶν ἕνα διαλειπόντων].}

{\greekfont Εἰ γὰρ δυνατόν, ἔστω ὁ Γ τετράγωνος. ἔστι δὲ καὶ ὁ Β τετράγωνος;
οἱ Β, Γ ἄρα πρὸς ἀλλήλους λόγον ἔθουσιν, ὃν τετράγωνος
ἀριθμὸς πρὸς τετράγωνον ἀριθμόν. καί ἐστιν ὡς ὁ Β πρὸς τὸν Γ,
ὁ Α πρὸς τὸν Β; οἱ Α, Β ἄρα πρὸς ἀλλήλους λόγον ἔθουσιν, ὃν τετράγωνος ἀριθμὸς πρὸς τετράγωνον ἀριθμόν; ὥστε οἱ Α, Β
ὅμοιοι ἐπίπεδοί εἰσιν. καί ἐστι τετράγωνος ὁ Β; τετράγωνος
ἄρα ἐστὶ καὶ ὁ Α; ὅπερ οὐθ ὑπέκειτο. οὐκ ἄρα ὁ Γ τετράγωνός
ἐστιν. ὁμοίως δὴ δείχομεν, ὅτι οὐδ ἄλλος οὐδεὶς τετράγωνός
ἐστι θωρὶς τοῦ τρίτου ἀπὸ τῆς μονάδος  καὶ τῶν ἕνα διαλειπόντων.}

{\greekfont Ἀλλὰ δὴ μὴ ἔστω ὁ Α κύβος. λέγω, ὅτι οὐδ ἄλλος οὐδεὶς
κύβος ἔσται θωρὶς τοῦ τετάρτου ἀπὸ τῆς μονάδος καὶ τῶν δύο διαλειπόντων.}

{\greekfont Εἰ γὰρ δυνατόν, ἔστω ὁ Δ κύβος. ἔστι δὲ καὶ ὁ Γ κύβος;
τέταρτος γάρ ἐστιν ἀπὸ τῆς μονάδος. καί ἐστιν ὡς ὁ Γ πρὸς
τὸν Δ, ὁ Β πρὸς τὸν Γ; καὶ ὁ Β ἄρα πρὸς τὸν Γ λόγον ἔχει, ὃν
κύβος πρὸς κύβον. καί ἐστιν ὁ Γ κύβος; καὶ ὁ Β ἄρα κύβος
ἐστίν. καὶ ἐπεί ἐστιν ὡς ἡ μονὰς  πρὸς τὸν Α, ὁ Α πρὸς
τὸν Β, ἡ δὲ μονὰς τὸν Α μετρεῖ κατὰ τὰς ἐν α ὐτῷ μονάδας,
καὶ ὁ Α ἄρα τὸν Β μετρεῖ κατὰ τὰς ἐν α ὐτῷ μονάδας;
ὁ Α ἄρα ἑαυτὸν πολλαπλασιάσας κύβον τὸν Β πεποίηκεν. ἐὰν
δὲ ἀριθμὸς ἑαυτὸν πολλαπλασιάσας κύβον ποιῇ, καὶ α ὐτὸς κύβος
ἔσται. κύβος ἄρα καὶ ὁ Α; ὅπερ οὐθ ὑπόκειται. οὐκ ἄρα
ὁ Δ κύβος ἐστίν. ὁμοίως δὴ δείχομεν, ὅτι οὐδ
ἄλλος οὐδεὶς κύβος ἐστὶ θωρὶς τοῦ τετάρτου ἀπὸ τῆς
μονάδος καὶ τῶν δύο διαλειπόντων; ὅπερ ἔδει δεῖχαι.}}

\ParallelRText{
\begin{center}
{\large Proposition 10}
\end{center}

If any multitude whatsoever of numbers is [continuously] proportional, (starting) from a unit, and the (number) after the
unit is not  square, then no other (number) will be  square
either, apart from the third from the unit, and all those (numbers after that) which leave an interval of  one (number). And if the (number) after the unit is not cube, then no other (number) will be cube
 either, apart from the fourth from the unit, and all those (numbers after that) which leave an interval of  two (numbers).
 
\epsfysize=1.8in
\centerline{\epsffile{Book09/fig08e.eps}}
 
Let any multitude whatsoever of numbers, $A$, $B$, $C$, $D$, $E$, $F$, be continuously proportional, (starting) from a unit. And let the (number) after the unit, $A$, not be  square. I say that  no other (number)  will
be square either, apart from the third from the unit [and (all) those (numbers after that) which leave an interval of  one (number)].

For, if possible, let $C$ be square.  And $B$ is also square [Prop.~9.8]. Thus, $B$ and $C$ have to one another (the) ratio which (some) square number (has) to (some other) square number. And as $B$ is to $C$, (so) $A$ (is) to $B$. Thus, $A$ and $B$ have
to one another (the) ratio which (some) square number has to (some other)
square number.  Hence, $A$ and $B$ are similar plane (numbers) [Prop.~8.26]. And $B$ is square.
Thus, $A$ is also  square. The very opposite thing was assumed.
$C$ is thus not square. So, similarly, we can show that no other
(number is) square either, apart from the  third from the unit, and (all) those (numbers after that) which leave an interval of  one (number).

And so let $A$ not be  cube. I say that  no other (number)  will
be cube either, apart from the fourth from the unit, and (all) those (numbers after that) which leave an interval of  two (numbers).

For, if possible, let $D$ be  cube. And $C$ is also  cube 
[Prop.~9.8]. For it is the fourth (number) from the unit.
And as $C$ is to $D$,  (so) $B$ (is) to $C$. And $B$ thus has to $C$ the
ratio which (some) cube (number has) to (some other) cube (number). And $C$
is  cube. Thus, $B$ is also cube [Props.~7.13, 8.25].
And since as the unit is to $A$,  (so) $A$ (is) to $B$, and the unit measures
$A$ according to the units in it, $A$ thus also measures $B$ according to the units in ($A$). Thus, $A$ has made the cube  (number) $B$ (by) multiplying
itself. And if a number makes a cube (number by) multiplying itself then
it itself will be  cube  [Prop.~9.6]. Thus, $A$
(is) also  cube. The very opposite thing was assumed. Thus, $D$ is not
cube. So, similarly, we can show that no other (number) is cube
either, apart from the fourth from the unit, and (all) those (numbers after that) which leave an interval of  two (numbers). (Which is) the very thing it was required to show.}
\end{Parallel}

%%%%%%
% Prop 9.11
%%%%%%
\pdfbookmark[1]{Proposition 9.11}{pdf9.11}
\begin{Parallel}{}{} 
\ParallelLText{
\begin{center}
{\large \ggn{11}.}
\end{center}\vspace*{-7pt}

{\greekfont Ἐὰν ἀπὸ μονάδος ὁποσοιοῦν ἀριθμοὶ ἑχῆς ἀνάλογον ὦσιν, ὁ ἐλάττων τὸν μείζονα μετρεῖ κατά τινα τῶν ὑπαρθόντ-ων ἐν τοῖς
ἀνάλογον ἀριθμοῖς.}\\

\epsfysize=1.6in
\centerline{\epsffile{Book09/fig11g.eps}}

{\greekfont Ἔστωσαν ἀπὸ μονάδος τῆς Α ὁποσοιοῦν ἀριθμοὶ ἑχῆς
ἀνάλογον οἱ Β, Γ, Δ, Ε; λέγω, ὅτι τῶν Β, Γ, Δ, Ε ὁ ἐλάθιστος
ὁ Β τὸν Ε μετρεῖ κατά τινα τῶν Γ, Δ.}

{\greekfont Ἐπεὶ γάρ ἐστιν ὡς ἡ Α μονὰς πρὸς τὸν Β, οὕτως ὁ Δ πρὸς
τὸν Ε, ἰσάκις ἄρα ἡ Α μονὰς τὸν Β ἀριθμὸν μετρεῖ καὶ ὁ Δ τὸν 
Ε; ἐναλλὰχ ἄρα ἰσάκις ἡ Α μονὰς τὸν Δ μετρεῖ καὶ ὁ Β τὸν Ε.
ἡ δὲ Α μονὰς τὸν Δ μετρεῖ  κατὰ τὰς ἐν α ὐτῷ μονάδας; καὶ ὁ Β ἄρα τὸν Ε μετρεῖ κατὰ τὰς ἐν τῷ Δ μονάδας; ὥστε ὁ ἐλάσσων ὁ Β τὸν μείζονα τὸν Ε μετρεῖ κατά
τινα ἀριθμὸν τῶν ὑπαρθόντων ἐν τοῖς ἀνάλογον ἀριθμοῖς.}\\~\\~\\

\begin{center}
{\large {\greekfont Πόρισμα}.}
\end{center}\vspace*{-7pt}

{\greekfont Καὶ φανερόν, ὅτι ἣν ἔχει τάχιν ὁ μετρῶν ἀπὸ μονάδος, τὴν
α ὐτὴν ἔχει καὶ ὁ καθ ὃν μετρεῖ ἀπὸ τοῦ μετρουμένου
ἐπὶ τὸ πρὸ α ὐτοῦ. ὅπερ ἔδει δεῖχαι.}}

\ParallelRText{
\begin{center}
{\large Proposition 11}
\end{center}

If any multitude whatsoever of numbers is continuously proportional, (starting) from a unit, then a lesser (number)
measures a greater according to some existing (number)  among
the proportional numbers.

\epsfysize=1.6in
\centerline{\epsffile{Book09/fig11e.eps}}

Let any multitude whatsoever of numbers,  $B$, $C$, $D$, $E$, be continuously proportional, (starting) from the unit $A$. I say that, for $B$, $C$, $D$, $E$, the least (number), $B$, measures $E$ according to some (one) of
$C$, $D$.

For since as the unit $A$ is to $B$, so $D$ (is) to $E$,  the unit $A$ thus measures the number $B$ the same number of times as $D$ (measures) $E$. Thus,
alternately, the unit $A$ measures $D$ the same number of times as $B$
(measures) $E$ [Prop.~7.15]. And the unit $A$ measures $D$ according to the units in it. Thus, $B$ also measures $E$ according to the units in $D$. Hence, the lesser (number) $B$
measures the greater $E$ according to some existing number among the
proportional numbers (namely, $D$). \\

\begin{center}
{\large Corollary}
\end{center}\vspace*{-7pt}

And (it is) clear that what(ever relative) place the measuring (number) has from the unit,
the  (number) according to which it measures has the same (relative) place from the measured (number), in (the direction of the number)
before it. (Which is) the very thing it was required to show.}
\end{Parallel}

%%%%%%
% Prop 9.12
%%%%%%
\pdfbookmark[1]{Proposition 9.12}{pdf9.12}
\begin{Parallel}{}{} 
\ParallelLText{
\begin{center}
{\large \ggn{12}.}
\end{center}\vspace*{-7pt}

{\greekfont Ἐὰν ἀπὸ μονάδος ὁποσοιοῦν ἀριθμοὶ ἑχῆς ἀνάλογον ὦσιν,
ὑφ ὅσων ἂν ὁ ἔσθατος πρώτων ἀριθμῶν μετρῆται, ὑπὸ
τῶν α ὐτῶν καὶ ὁ παρὰ τὴν μονάδα μετρηθήσεται.}\\

\epsfysize=1.1in
\centerline{\epsffile{Book09/fig12g.eps}}

{\greekfont Ἔστωσαν ἀπὸ μονάδος ὁποσοιδηποτοῦν ἀριθμοὶ ἀνάλογ-ον οἱ 
Α, Β, Γ, Δ; λέγω, ὅτι ὑφ ὅσων ἂν ὁ Δ πρώτων ἀριθμῶν μετρῆται, ὑπὸ τῶν α ὐτῶν καὶ ὁ Α μετρηθήσεται.}

{\greekfont Μετρείσθω γὰρ ὁ Δ ὑπό τινος πρώτου ἀριθμοῦ τοῦ Ε; λέγω, ὅτι
ὁ Ε τὸν Α μετρεῖ. μὴ γάρ; καί ἐστιν ὁ Ε πρῶτος, ἅπας δὲ πρῶτος
ἀριθμὸς πρὸς ἅπαντα, ὃν μὴ μετρεῖ, πρῶτός ἐστιν; οἱ Ε, Α ἄρα
πρῶτοι πρὸς ἀλλήλους εἰσίν. καὶ ἐπεὶ ὁ Ε τὸν Δ μετρεῖ, μετρείτω
α ὐτὸν κατὰ τὸν Ζ; ὁ Ε ἄρα τὸν Ζ πολλαπλασιάσας τὸν Δ πεποίηκεν. πάλιν, ἐπεὶ ὁ Α τὸν Δ μετρεῖ κατὰ τὰς ἐν τῷ Γ μονάδας, ὁ 
Α  ἄρα τὸν Γ πολλαπλασιάσας τὸν Δ πεποίηκεν. ἀλλὰ μὴν καὶ ὁ Ε
τὸν Ζ πολλαπλασιάσας τὸν Δ πεποίηκεν; ὁ ἄρα ἐκ τῶν Α, Γ ἴσος
ἐστὶ τῷ ἐκ τῶν Ε, Ζ. ἔστιν ἄρα ὡς ὁ Α πρὸς τὸν Ε, ὁ Ζ
πρὸς τὸν Γ. οἱ δὲ Α, Ε πρῶτοι, οἱ δὲ πρῶτοι καὶ ἐλάθιστοι, οἱ
δὲ ἐλάθιστοι μετροῦσι τοὺς τὸν α ὐτὸν λόγον ἔθοντας ἰσάκις
ὅ τε ἡγούμενος τὸν ἡγούμενον καὶ ὁ ἑπόμενος τὸν
ἑπόμενον; μετρεῖ ἄρα ὁ Ε τὸν Γ. μετρείτω α ὐτὸν κατὰ τὸν Η;
ὁ Ε ἄρα τὸν Η πολλαπλασιάσας τὸν Γ πεποίηκεν. ἀλλὰ μὴν διὰ
τὸ πρὸ τούτου καὶ ὁ Α τὸν Β πολλαπλασιάσας τὸν Γ πεποίηκεν. ὁ
ἄρα ἐκ τῶν Α, Β ἴσος ἐστὶ τῷ ἐκ τῶν Ε, Η. ἔστιν ἄρα
ὡς ὁ Α πρὸς τὸν Ε, ὁ Η πρὸς τὸν Β. οἱ δὲ Α, Ε πρῶτοι, οἱ
δὲ πρῶτοι καὶ ἐλάθιστοι, οἱ δὲ ἐλάθιστοι ἀριθμοὶ μετροῦσι τοὺς
τὸν α ὐτὸν λόγον ἔθοντας α ὐτοῖς ἰσάκις ὅ τε ἡγούμενος
τὸν ἡγούμενον καὶ ὁ ἑπόμενος τὸν ἑπόμενον; μετρεῖ ἄρα ὁ Ε
τὸν Β. μετρείτω α ὐτὸν κατὰ τὸν Θ; ὁ Ε ἄρα τὸν Θ πολλαπλασιάσας
τὸν Β πεποίηκεν. ἀλλὰ μὴν καὶ ὁ Α ἑαυτὸν πολλαπλασιάσας τὸν Β
πεποίηκεν; ὁ ἄρα ἐκ τῶν Ε, Θ ἴσος ἐστὶ τῷ ἀπὸ τοῦ Α.
ἔστιν ἄρα ὡς ὁ Ε πρὸς τὸν Α, ὁ Α πρὸς τὸν Θ. οἱ δὲ Α, Ε
πρῶτοι, οἱ δὲ πρῶτοι καὶ ἐλάθιστοι, οἱ δὲ ἐλάθιστοι μετροῦσι
τοὺς τὸν α ὐτὸν λόγον ἔθοντας ἰσάκις ὅ ἡγούμενος τὸν ἡγούμενον καὶ ὁ ἑπόμενος τὸν ἑπόμενον; μετρεῖ ἄρα
ὁ Ε τὸν Α ὡς ἡγούμενος ἡγούμενον. ἀλλὰ μὴν καὶ οὐ
μετρεῖ; ὅπερ ἀδύνατον. οὐκ ἄρα οἱ Ε, Α πρῶτοι πρὸς ἀλλήλους
εἰσίν. σύνθετοι ἄρα. οἱ δὲ σύνθετοι ὑπὸ [πρώτου] ἀριθμοῦ τινος
μετροῦνται. καὶ ἐπεὶ ὁ Ε πρῶτος ὑπόκειται, ὁ δὲ πρῶτος
ὑπὸ ἑτέρου ἀριθμοῦ οὐ μετρεῖται ἢ ὑφ ἑαυτοῦ, ὁ Ε
ἄρα τοὺς Α, Ε μετρεῖ; ὥστε ὁ Ε τὸν Α μετρεῖ. μετρεῖ δὲ καὶ
τὸν Δ; ὁ Ε ἄρα τοὺς Α, Δ μετρεῖ. ὁμοίως δὴ δείχομεν, ὅτι
ὑφ ὅσων ἂν ὁ Δ πρώτων ἀριθμῶν μετρῆται, ὑπὸ τῶν α ὐτῶν
καὶ ὁ Α μετρηθήσεται; ὅπερ ἔδει δεῖχαι.}}

\ParallelRText{
\begin{center}
{\large Proposition 12}
\end{center}

If  any multitude whatsoever of numbers is continuously proportional, (starting) from a unit, then  however  many prime numbers  the
last (number) is measured by, the (number) next to the unit will also be measured by the same (prime numbers).

\epsfysize=1.1in
\centerline{\epsffile{Book09/fig12e.eps}}

Let any multitude whatsoever of numbers,  $A$, $B$, $C$, $D$, be (continuously) proportional, (starting) from a unit. I say that however many prime
numbers $D$ is measured by, $A$ will also be measured by the same
(prime numbers).

For let $D$ be measured by some prime number $E$. I say that $E$ measures $A$. For (suppose it does) not. $E$ is prime, and every prime number is prime to
every number which it does not measure [Prop.~7.29]. Thus, $E$ and $A$ are prime to one another. And since $E$ measures $D$, let it measure it according to $F$. 
Thus, $E$ has made $D$ (by) multiplying $F$.  Again, since $A$ measures $D$ according to the units in $C$ [Prop.~9.11~corr.],
$A$ has thus made $D$ (by) multiplying $C$. But, in fact, $E$ has also
made $D$ (by) multiplying $F$. Thus, the (number created) from (multiplying)
$A$, $C$ is equal to  the (number created) from (multiplying) $E$, $F$.
Thus, as $A$ is to $E$, (so) $F$ (is) to $C$ [Prop.~7.19]. And $A$ and $E$ (are)  prime (to one another), and (numbers) prime (to one another are) also the least (of those
numbers having the same ratio as them) [Prop.~7.21],
and the least (numbers) measure those (numbers) having the same ratio as them an equal number of times, the leading (measuring) the leading, and the
following the following [Prop.~7.20]. Thus, $E$
measures $C$. Let it measure it according to $G$. Thus, $E$ has made $C$
(by) multiplying $G$. But, in fact, via the (proposition) before this, $A$
has also made $C$ (by) multiplying $B$ [Prop.~9.11~corr.].
 Thus, the (number created) from (multiplying) $A$, $B$ is equal to the
 (number created) from (multiplying) $E$, $G$. Thus, as $A$ is to $E$, (so) $G$ (is) to $B$ [Prop.~7.19]. And $A$ and $E$
 (are) prime (to one another), and (numbers) prime (to one another are) also the least (of those
numbers having the same ratio as them) [Prop.~7.21],
and the least (numbers) measure those (numbers) having the same ratio as them an equal number of times, the leading (measuring) the leading, and the
following the following [Prop.~7.20]. Thus, $E$
measures $B$. Let it measure it according to $H$. Thus, $E$ has made $B$ (by) multiplying $H$. But, in fact, $A$ has also made $B$
(by) multiplying itself [Prop.~9.8].  Thus, the (number created) from (multiplying) $E$, $H$ is equal to the (square) on $A$.
Thus, as $E$ is to $A$, (so) $A$ (is) to $H$ [Prop.~7.19]. And $A$ and $E$
 are prime (to one another), and (numbers) prime (to one another are) also the least (of those
numbers having the same ratio as them) [Prop.~7.21],
and the least (numbers) measure those (numbers) having the same ratio  as them an equal number of times, the leading (measuring) the leading, and the
following the following
[Prop.~7.20]. Thus, $E$
measures $A$, as the leading (measuring the) leading. But, in fact, 
($E$) also does not measure ($A$). The very thing (is) impossible. Thus,
$E$ and $A$ are not prime to one another. Thus, (they are) composite (to one another). And (numbers) composite (to one another) are (both) measured
by some [prime] number [Def.~7.14]. And since $E$
is assumed (to be) prime, and a prime (number) is not measured by  another number (other) than itself
[Def.~7.11], $E$ thus measures (both) $A$ and $E$. Hence, $E$ measures $A$. And it also measures $D$. Thus, $E$ measures (both) $A$ and $D$. So, similarly, we can show that however many prime
numbers $D$ is measured by, $A$ will also be measured by the same
(prime numbers). (Which is) the very thing it was required to show.}
\end{Parallel}

%%%%%%
% Prop 9.13
%%%%%%
\pdfbookmark[1]{Proposition 9.13}{pdf9.13}
\begin{Parallel}{}{} 
\ParallelLText{
\begin{center}
{\large \ggn{13}.}
\end{center}\vspace*{-7pt}

{\greekfont Ἐὰν ἀπὸ μονάδος ὁποσοιοῦν ἀριθμοὶ ἑχῆς ἀνάλογον ὦσιν,
ὁ δὲ μετὰ τὴν μονάδα πρῶτος ᾖ, ὁ μέγιστος ὑπ οὐδενὸς 
[ἄλλου] μετρηθήσεται παρὲχ τῶν ὑπαρθόντων  ἐν τοῖς ἀνάλογον
ἀριθμοῖς.}

{\greekfont Ἔστωσαν ἀπὸ μονάδος ὁποσοιοῦν ἀριθμοὶ ἑχῆς ἀνάλογον οἱ 
Α, Β, Γ, Δ, ὁ δὲ μετὰ τὴν μονάδα ὁ Α πρῶτος ἔστω; λέγω, ὅτι
ὁ μέγιστος α ὐτῶν ὁ Δ ὑπ οὐδενὸς ἄλλου μετρηθήσεται
παρὲχ τῶν Α, Β, Γ.}\\~\\

\epsfysize=1.1in
\centerline{\epsffile{Book09/fig13g.eps}}

{\greekfont Εἰ γὰρ δυνατόν, μετρείσθω ὑπὸ τοῦ Ε, καὶ ὁ Ε μ ηδενὶ τῶν Α, Β, Γ
ἔστω ὁ α ὐτός. φανερὸν δή, ὅτι ὁ Ε πρῶτος οὔκ ἐστιν. εἰ 
γὰρ ὁ Ε πρῶτός ἐστι καὶ μετρεῖ τὸν Δ, καὶ τὸν Α μετρήσει
πρῶτον ὄντα μὴ ὢν α ὐτῷ ὁ α ὐτός; ὅπερ ἐστὶν ἀδύνατον.
οὐκ ἄρα ὁ Ε πρῶτός ἐστιν. σύνθετος ἄρα. πᾶς δὲ σύνθετος ἀριθμὸς
ὑπὸ πρώτου τινὸς ἀριθμοῦ μετρεῖται; ὁ Ε ἄρα ὑπὸ πρώτου τινὸς
ἀριθμοῦ μετρεῖται.
λέγω δή, ὅτι ὑπ οὐδενὸς
ἄλλου πρώτου μετρηθήσεται πλὴν τοῦ Α. εἰ γὰρ ὑφ ἑτέρου μετρεῖται
ὁ Ε, ὁ δὲ Ε τὸν Δ μετρεῖ, κἀκεῖνος ἄρα τὸν Δ μετρήσει; ὥστε καὶ
τὸν Α  μετρήσει πρῶτον ὄντα μὴ ὢν α ὐτῷ ὁ α ὐτός; ὅπερ
ἐστὶν ἀδύνατον. ὁ Α ἄρα τὸν  Ε μετρεῖ. καὶ ἐπεὶ ὁ Ε τὸν Δ
μετρεῖ, μετρείτω α ὐτὸν κατὰ τὸν Ζ. λέγω, ὅτι ὁ Ζ οὐδενὶ τῶν
Α, Β, Γ ἐστιν ὁ α ὐτός. εἰ γὰρ ὁ Ζ ἑνὶ τῶν Α, Β, Γ ἐστιν ὁ
α ὐτὸς καὶ μετρεῖ τὸν Δ κατὰ τὸν Ε, καὶ εἷς ἄρα τῶν Α, Β, Γ
τὸν Δ μετρεῖ κατά τὸν Ε. ἀλλὰ εῖ̔ς τῶν Α, Β, Γ τὸν Δ μετρεῖ κατά
τινα τῶν Α, Β, Γ; καὶ ὁ Ε ἄρα ἑνὶ τῶν Α, Β, Γ
ἐστιν ὁ α ὐτός; ὅπερ οὐθ ὑπόκειται. οὐκ ἄρα ὁ Ζ ἑνὶ τῶν
Α, Β, Γ ἐστιν ὁ α ὐτός. ὁμοίως δὴ δείχομεν, ὅτι μετρεῖται ὁ Ζ
ὑπὸ τοῦ Α, δεικνύντες πάλιν, ὅτι ὁ Ζ οὔκ ἐστι πρῶτος. εἰ γὰρ,
καὶ μετρεῖ τὸν Δ, καὶ τὸν Α μετρήσει πρῶτον ὄντα μὴ ὢν α ὐτῷ
ὁ α ὐτός; ὅπερ ἐστὶν ἀδύνατον; οὐκ ἄρα πρῶτός ἐστιν ὁ Ζ;
σύνθετος ἄρα. ἅπας δὲ σύνθετος ἀριθμὸς ὑπὸ πρώτου τινὸς ἀριθμοῦ
μετρεῖται; ὁ Ζ ἄρα ὑπὸ πρώτου τινὸς ἀριθμοῦ μετρεῖται. λέγω δή,
ὅτι ὑφ ἑτέρου πρώτου οὐ μετρηθήσεται πλὴν τοῦ Α. εἰ γὰρ ἕτερός
τις πρῶτος τὸν Ζ μετρεῖ, ὁ δὲ Ζ τὸν Δ μετρεῖ, κἀκεῖνος ἄρα τὸν Δ
μετρήσει; ὥστε καὶ τὸν Α μετρήσει πρῶτον ὄντα μὴ ὢν α ὐτῷ
ὁ α ὐτός; ὅπερ ἐστὶν ἀδύνατον. ὁ Α ἄρα τὸν Ζ μετρεῖ. καὶ ἐπεὶ
ὁ Ε τὸν Δ μετρεῖ κατὰ τὸν Ζ, ὁ Ε ἄρα τὸν Ζ πολλαπλασιάσας
τὸν Δ πεποίηκεν. ἀλλὰ μὴν καὶ ὁ Α τὸν Γ πολλαπλασιάσας
τὸν Δ πεποίηκεν; ὁ ἄρα ἐκ τῶν Α, Γ ἴσος ἐστὶ τῷ ἐκ τῶν 
Ε, Ζ. ἀνάλογον ἄρα ἐστὶν ὡς ὁ Α πρὸς τὸν Ε, οὕτως ὁ Ζ
πρὸς τὸν Γ. ὁ δὲ Α τὸν Ε μετρεῖ; καὶ ὁ Ζ ἄρα τὸν Γ μετρεῖ. μετρείτω α ὐτὸν κατὰ τὸν Η. ὁμοίως δὴ δείχομεν, ὅτι 
ὁ Η οὐδενὶ τῶν Α, Β ἐστιν ὁ α ὐτός, καὶ ὅτι μετρεῖται ὑπὸ
τοῦ Α. καὶ ἐπεὶ ὁ Ζ τὸν Γ μετρεῖ κατὰ τὸν Η,  ὁ Ζ ἄρα τὸν Η
πολλαπλασιάσας τὸν Γ πεποίηκεν. ἀλλὰ μὴν καὶ ὁ Α τὸν Β πολλαπλασιάσας
τὸν Γ πεποίηκεν; ὁ ἄρα ἐκ τῶν Α, Β ἴσος ἐστὶ τῷ ἐκ τῶν Ζ, Η.
ἀνάλογον ἄρα ὡς ὁ Α πρὸς τὸν Ζ, ὁ Η πρὸς τὸν  Β.  μετρεῖ δὲ ὁ
Α τὸν Ζ; μετρεῖ ἄρα καὶ ὁ Η τὸν Β. μετρείτω
α ὐτὸν κατὰ τὸν Θ. ὁμοίως δὴ
δείχομεν, ὅτι ὁ Θ τῷ Α οὐκ ἔστιν
ὁ α ὐτός. καὶ ἐπεὶ ὁ Η τὸν Β μετρεῖ κατὰ τὸν Θ,  ὁ Η ἄρα τὸν
Θ πολλαπλασιάσας τὸν Β πεποίηκεν. ἀλλὰ μὴν καὶ ὁ Α ἑαυτὸν πολλαπλασιάσας τὸν Β πεποίηκεν; ὁ ἄρα ὑπὸ Θ, Η ἴσος ἐστὶ τῷ
ἀπὸ τοῦ Α τετραγώνῳ;  ἔστιν ἄρα ὡς ὁ Θ πρὸς τὸν Α, ὁ Α
πρὸς τὸν Η. μετρεῖ δὲ ὁ Α τὸν Η; μετρεῖ ἄρα καὶ ὁ Θ τὸν Α
πρῶτον ὄντα μὴ ὢν α ὐτῷ ὁ α ὐτός; ὅπερ ἄτοπον.
οὐκ ἄρα ὁ μέγιστος ὁ Δ ὑπὸ ἑτέρου ἀριθμοῦ μετρηθήσεται
παρὲχ τῶν Α, Β, Γ; ὅπερ ἔδει δεῖχαι.}}

\ParallelRText{
\begin{center}
{\large Proposition 13}
\end{center}

If  any multitude whatsoever of numbers is continuously proportional, (starting) from a unit, and the (number) after the
unit is prime, then the greatest (number) will be measured by no [other] (numbers) except (numbers) existing  among the proportional numbers.

Let any multitude whatsoever of numbers,  $A$, $B$, $C$, $D$, be continuously proportional, (starting) from a unit. And let the (number)
after the unit, $A$, be prime. I say that the greatest of them, $D$, will be
measured by no other (numbers) except $A$, $B$, $C$. 

\epsfysize=1.1in
\centerline{\epsffile{Book09/fig13e.eps}}

For, if possible, let it be measured by $E$, and let $E$ not be
the same as one of $A$, $B$, $C$. So it is clear that $E$ is not prime. 
For if $E$ is prime, and measures $D$, then it will also measure $A$, (despite $A$) being prime (and) not being  the same as it [Prop.~9.12]. The very thing is impossible. Thus, $E$
is not prime. Thus, (it is) composite. And every composite
number is measured by some prime number [Prop.~7.31]. Thus, $E$ is measured by some prime number. So I say that it will be measured by no other prime number than $A$.
For if $E$ is measured by another (prime number), and $E$ measures $D$, then this (prime number) will thus also measure $D$. Hence, it will also
measure $A$,  (despite $A$) being prime (and) not being  the same as it [Prop.~9.12]. The very thing is impossible. Thus,
$A$ measures $E$. And since $E$ measures $D$, let it measure it according to $F$. I say that $F$ is not the same as one of $A$, $B$, $C$. For if $F$ is
the same as one of $A$, $B$, $C$, and measures $D$ according to $E$,
then one of $A$, $B$, $C$ thus also measures $D$ according to $E$.
But one of $A$, $B$, $C$ (only) measures $D$ according to some (one) of
$A$, $B$, $C$
[Prop.~9.11]. And thus $E$ is the same as one of
$A$, $B$, $C$. The very opposite thing was assumed. Thus, $F$ is not the
same as one of $A$, $B$, $C$. Similarly, we can show that $F$ is measured by $A$,  (by) again showing that $F$ is not prime. For if ($F$ is prime), and measures $D$,  then it will also measure $A$, (despite $A$) being prime (and) not being  the same as it [Prop.~9.12]. The very thing
is impossible. Thus, $F$ is not prime. Thus, (it is) composite.  And every composite
number is measured by some prime number [Prop.~7.31]. Thus, $F$ is measured by some prime
number. So I say that it will be measured by no other prime number than $A$.
For if  some other prime (number) measures $F$, and $F$ measures $D$, then this (prime number) will thus also measure $D$.  Hence, it will also
measure $A$,  (despite $A$) being prime (and) not being  the same as it [Prop.~9.12]. The very thing is impossible.
Thus, $A$ measures $F$. And since $E$ measures $D$ according to 
$F$, $E$ has thus made $D$ (by) multiplying $F$.  But, in fact, $A$ has
also made $D$ (by) multiplying $C$ [Prop.~9.11~corr.]. Thus, the (number created) from (multiplying)
$A$, $C$ is equal to  the (number created) from (multiplying) $E$, $F$. Thus, proportionally, as $A$ is to $E$, so $F$ (is) to $C$ [Prop.~7.19]. And $A$ measures $E$. Thus, $F$ also
measures $C$. Let it
 measure it according to $G$. So, similarly, we can show
that $G$ is not the same as one of $A$, $B$, and that it is measured by $A$. 
And since $F$ measures $C$ according to $G$, $F$ has thus made
$C$ (by) multiplying $G$.  But, in fact, $A$ has
also made $C$ (by) multiplying $B$ [Prop.~9.11~corr.]. Thus, the (number created) from (multiplying)
$A$, $B$ is equal to  the (number created) from (multiplying) $F$, $G$.
Thus, proportionally, as $A$ (is) to $F$, so $G$ (is) to $B$ [Prop.~7.19]. And $A$ measures $F$. 
Thus, $G$ also measures $B$.
Let it measure
it according to $H$. So, similarly, we can show that $H$ is not the same as $A$. And since $G$ measures $B$ according to $H$, $G$ has thus made $B$
(by) multiplying $H$.  But, in fact, $A$ has
 also made $B$ (by) multiplying itself [Prop.~9.8]. Thus, the (number created) from (multiplying)
$H$, $G$ is equal to  the square on $A$. Thus, as $H$ is to $A$, (so) $A$ (is) to $G$ [Prop.~7.19]. And $A$ measures $G$.
Thus, $H$ also measures $A$, (despite $A$) being prime (and) not being  the same as it. The very thing (is) absurd. Thus, the greatest (number) $D$ cannot be
measured by  another (number) except (one of) $A$, $B$, $C$. (Which is) the
very thing it was required to show.}
\end{Parallel}

%%%%%%
% Prop 9.14
%%%%%%
\pdfbookmark[1]{Proposition 9.14}{pdf9.14}
\begin{Parallel}{}{} 
\ParallelLText{
\begin{center}
{\large \ggn{14}.}
\end{center}\vspace*{-7pt}

{\greekfont Ἐὰν ἐλάθιστος ἀριθμὸς ὑπὸ πρώτων ἀριθμῶν μετρῆται, ὑπ
οὑδενὸς ἄλλου πρώτου ἀριθμοῦ μετρηθήσεται παρὲχ τῶν ἐχ
ἀρθῆς μετρούντων.}

{\greekfont Ἐλάθιστος γὰρ ἀριθμὸς ὁ Α ὑπὸ πρώτων ἀριθμῶν τῶν Β, Γ, Δ
μετρείσθω; λέγω, ὅτι ὁ Α ὑπ οὐδενὸς ἄλλου πρώτου
ἀριθμοῦ μετρηθήσεται παρὲχ τῶν Β, Γ, Δ.}

\epsfysize=0.8in
\centerline{\epsffile{Book09/fig14g.eps}}

{\greekfont Εἰ γὰρ δυνατόν, μετρείσθω ὑπὸ πρώτου τοῦ Ε, καὶ ὁ Ε μ ηδενὶ
τῶν Β, Γ, Δ ἔστω ὁ α ὐτός. καὶ ἐπεὶ ὁ Ε τὸν Α μετρεῖ, μετρείτω
α ὐτὸν κατὰ τὸν Ζ; ὁ Ε ἄρα τὸν Ζ πολλαπλασιάσας τὸν Α πεποίηκεν.
καὶ μετρεῖται ὁ Α ὑπὸ πρώτων ἀριθμῶν τῶν Β, Γ, Δ. ἐὰν
δὲ δύο ἀριθμοὶ πολλαπλασιάσαντες ἀλλήλους ποιῶσί τινα, τὸν δὲ
γενόμενον ἐχ α ὐτῶν μετρῇ τις πρῶτος ἀριθμός, καὶ
ἕνα τῶν ἐχ ἀρθῆς μετρήσει; οἱ Β, Γ, Δ ἄρα ἕνα τῶν Ε, Ζ
μετρήσουσιν. τὸν  μὲν οὖν Ε οὐ μετρήσουσιν; ὁ γὰρ Ε πρῶτός
ἐστι καὶ οὐδενὶ τῶν Β, Γ, Δ ὁ α ὐτός. τὸν Ζ ἄρα μετροῦσιν
ἐλάσσονα ὄντα τοῦ Α; ὅπερ ἀδύνατον. ὁ γὰρ Α ὑπόκειται
ἐλάθιστος ὑπὸ τῶν Β, Γ, Δ μετρούμενος. οὐκ ἄρα τὸν
Α μετρήσει πρῶτος ἀριθμὸς παρὲχ τῶν Β, Γ, Δ; ὅπερ
ἔδει δεῖχαι.}}

\ParallelRText{
\begin{center}
{\large Proposition 14}
\end{center}

If  a least number is measured by (some) prime numbers then it will not be measured by
any other prime number except (one of) the original measuring (numbers).

For let $A$ be the least number measured by the prime numbers $B$, $C$, $D$. I say that $A$ will not be measured by any other prime number except
(one of) $B$, $C$, $D$.

\epsfysize=0.9in
\centerline{\epsffile{Book09/fig14e.eps}}

For, if possible, let it be measured by the prime (number) $E$. And
let $E$ not be the same as one of $B$, $C$, $D$. And since $E$ measures
$A$, let it measure it according to $F$.  Thus, $E$ has made $A$ (by)
multiplying $F$. And $A$ is measured by the prime numbers $B$, $C$, $D$.
And if two numbers make some (number by) multiplying one another,
and some prime number  measures the number created from them, then (the prime number)
will also measure one of the original (numbers) [Prop.~7.30]. Thus, $B$, $C$, $D$ will measure one of
$E$, $F$. In fact, they do not measure $E$. For $E$ is prime, and not the
same as one of $B$, $C$, $D$. Thus, they (all) measure $F$, which is less than $A$. The very thing (is) impossible. For $A$ was assumed
(to be) the least (number) measured by $B$, $C$, $D$. Thus, no prime number can measure $A$ except (one of) $B$, $C$, $D$. (Which is) the very thing it was required to show.}
\end{Parallel}

%%%%%%
% Prop 9.15
%%%%%%
\pdfbookmark[1]{Proposition 9.15}{pdf9.15}
\begin{Parallel}{}{} 
\ParallelLText{
\begin{center}
{\large \ggn{15}.}
\end{center}\vspace*{-7pt}

{\greekfont Ἐὰν τρεῖς ἀριθμοὶ ἑχῆς ἀνάλογον ὦσιν ἐλάθιστοι τῶν τὸν
α ὐτὸν λόγον ἐθόντων α ὐτοῖς, δύο ὁποιοιοῦν συντεθέντες πρὸς
τὸν λοιπὸν πρῶτοί εἰσιν.}\\

\epsfysize=1.in
\centerline{\epsffile{Book09/fig15g.eps}}

{\greekfont Ἔστωσαν τρεῖς ἀριθμοὶ ἑχῆς ἀνάλογον ἐλάθιστοι τῶν τὸν α ὐτὸν
λόγον ἐθόντων α ὐτοῖς οἱ Α, Β, Γ;
λέγω, ὅτι τῶν Α, Β, Γ δύο ὁποιοιοῦν συντεθέντες
πρὸς τὸν λοιπὸν πρῶτοι εἰσιν, οἱ μὲν Α, Β πρὸς τὸν Γ, οἱ δὲ
Β, Γ πρὸς τὸν Α καὶ ἔτι οἱ Α, Γ πρὸς τὸν Β.}

{\greekfont Εἰλήφθωσαν γὰρ ἐλάθιστοι ἀριθμοὶ τῶν τὸν α ὐτὸν λόγον
ἐθόντων τοῖς Α, Β, Γ δύο οἱ ΔΕ, ΕΖ. φανερὸν δή, ὅτι ὁ μὲν
ΔΕ ἑαυτὸν πολλαπλασιάσας τὸν Α πεποίηκεν, τὸν δὲ ΕΖ πολλαπλασιάσας
τὸν Β πεποίηκεν, καὶ ἔτι ὁ ΕΖ ἑαυτὸν πολλαπλασιάσας τὸν Γ πεποίηκεν.
καὶ ἐπεὶ οἱ ΔΕ, ΕΖ ἐλάθιστοί εἰσιν, πρῶτοι πρὸς ἀλλήλους εἰσίν.
ἐὰν δὲ δύο ἀριθμοὶ πρῶτοι πρὸς ἀλλήλους ὦσιν, καὶ
συναμφότερος πρὸς ἑκάτερον πρῶτός ἐστιν; καὶ ὁ ΔΖ ἄρα
πρὸς ἑκάτερον τῶν ΔΕ, ΕΖ πρῶτός ἐστιν. ἀλλὰ μὴν καὶ ὁ ΔΕ
πρὸς τὸν ΕΖ πρῶτός ἐστιν; οἱ ΔΖ, ΔΕ ἄρα πρὸς τὸν ΕΖ πρῶτοί
εἰσιν. ἐὰν δὲ δύο ἀριθμοὶ πρός τινα ἀριθμὸν πρῶτοι ὦσιν, καὶ
ὁ ἐχ α ὐτῶν γενόμενος πρὸς τὸν λοιπὸν πρῶτός ἐστιν; ὥστε
ὁ ἐκ τῶν ΖΔ, ΔΕ πρὸς τὸν ΕΖ πρῶτός ἐστιν; ὥστε καὶ ὁ ἐκ
τῶν ΖΔ, ΔΕ πρὸς τὸν ἀπὸ τοῦ ΕΖ πρῶτός ἐστιν. [ἐὰν
γὰρ δύο ἀριθμοὶ πρῶτοι πρὸς ἀλλήλους ὦσιν, ὁ ἐκ τοῦ ἑνὸς
α ὐτῶν γενόμενος πρὸς τὸν λοιπὸν πρῶτός ἐστιν]. ἀλλ ὁ ἐκ
τῶν ΖΔ, ΔΕ ὁ ἀπὸ τοῦ ΔΕ ἐστι μετὰ τοῦ ἐκ τῶν ΔΕ, ΕΖ; ὁ
ἄρα ἀπὸ τοῦ ΔΕ μετὰ τοῦ ἐκ τῶν ΔΕ, ΕΖ πρὸς τὸν ἀπὸ τοῦ
 ΕΖ πρῶτός
ἐστιν. καί ἐστιν ὁ μὲν ἀπὸ τοῦ ΔΕ ὁ Α, ὁ δὲ ἐκ τῶν ΔΕ, ΕΖ
ὁ Β, ὁ δὲ ἀπὸ τοῦ ΕΖ ὁ Γ; οἱ Α, Β ἄρα συντεθέντες πρὸς τὸν Γ
πρῶτοί εἰσιν. ὁμοίως δὴ δείχομεν, ὅτι καὶ οἱ Β, Γ
πρὸς τὸν Α πρῶτοί εἰσιν. λέγω δή, ὅτι καὶ οἱ Α, Γ πρὸς τὸν Β πρῶτοί εἰσιν.
ἐπεὶ γὰρ ὁ ΔΖ πρὸς ἑκάτερον
τῶν ΔΕ, ΕΖ πρῶτός ἐστιν, καὶ ὁ ἀπὸ τοῦ ΔΖ πρὸς τὸν ἐκ τῶν
ΔΕ, ΕΖ πρῶτός ἐστιν. ἀλλὰ τῷ ἀπὸ τοῦ ΔΖ ἴσοι
εἰσὶν οἱ ἀπὸ τῶν ΔΕ, ΕΖ μετὰ τοῦ δὶς ἐκ τῶν ΔΕ, ΕΖ; καὶ
οἱ ἀπὸ τῶν ΔΕ, ΕΖ ἄρα μετὰ τοῦ δὶς ὑπὸ τῶν ΔΕ, ΕΖ πρὸς
τὸν ὑπὸ τῶν ΔΕ, ΕΖ πρῶτοί [εἰσι]. διελόντι οἱ ἀπὸ τῶν ΔΕ, ΕΖ
μετὰ τοῦ ἅπαχ ὑπὸ ΔΕ, ΕΖ πρὸς τὸν ὑπὸ ΔΕ, ΕΖ πρῶτοί
εἰσιν. ἔτι διελόντι οἱ ἀπὸ τῶν ΔΕ, ΕΖ ἄρα πρὸς τὸν 
ὑπὸ ΔΕ, ΕΖ πρῶτοί εἰσιν. καί ἐστιν ὁ μὲν ἀπὸ τοῦ ΔΕ ὁ
Α, ὁ δὲ ὑπὸ τῶν ΔΕ, ΕΖ ὁ Β, ὁ δὲ ἀπὸ τοῦ ΕΖ ὁ Γ. οἱ
Α, Γ ἄρα συντεθέντες πρὸς τὸν Β πρῶτοί εἰσιν; ὅπερ ἔδει
δεῖχαι.}}

\ParallelRText{
\begin{center}
{\large Proposition 15}
\end{center}

If three continuously proportional
numbers are the least of those (numbers) having the same ratio as them then
two (of them) added together in any way are prime to the remaining (one).

\epsfysize=1.in
\centerline{\epsffile{Book09/fig15e.eps}}

Let $A$, $B$, $C$ be three continuously proportional numbers (which are)
the least of those (numbers) having the same ratio as them. I say that two
of $A$, $B$, $C$ added together in any way are prime to the remaining
(one), (that is) $A$ and $B$ (prime) to $C$, $B$ and $C$ to
$A$, and, further, $A$ and $C$ to $B$.

Let the two least numbers, $DE$ and $EF$, having the same ratio as
$A$, $B$, $C$, be taken [Prop.~8.2]. So
it is clear that $DE$ has made $A$ (by) multiplying itself, and
has made $B$ (by) multiplying $EF$, and, further, $EF$ has made $C$ (by)
multiplying itself [Prop.~8.2]. And since $DE$, $EF$ are the least (of those numbers having the same ratio as them), 
they are prime to one another [Prop.~7.22]. 
And if two numbers are prime to one another then the sum (of them)
is also prime to each [Prop.~7.28]. Thus, $DF$
is also prime to each of $DE$, $EF$. But, in fact, $DE$ is  also prime to $EF$.
Thus, $DF$, $DE$ are (both) prime to $EF$. And if two numbers are (both)
prime to some number then the (number) created from (multiplying) them
is also prime to the remaining (number) [Prop.~7.24].
Hence, the (number created) from (multiplying) $FD$, $DE$ is prime
to $EF$. Hence, the (number created) from (multiplying) $FD$, $DE$ is also prime
to the (square) on $EF$ [Prop.~7.25]. 
[For if two numbers are prime to one another then the (number) created from
(squaring) one of them is prime to the remaining (number).]
But the (number created) from (multiplying) $FD$, $DE$ is the
(square) on $DE$ plus the (number created) from (multiplying) $DE$, $EF$
[Prop.~2.3]. Thus, the (square) on $DE$ plus
the (number created) from (multiplying) $DE$, $EF$ is prime to the
(square) on $EF$. 
And the (square)
on $DE$ is $A$, and the (number created) from (multiplying) $DE$, $EF$
(is) $B$, and the (square) on $EF$ (is) $C$. Thus,  $A$, $B$ summed is prime to $C$. So,  similarly, we can show that $B$, $C$ (summed) is also
prime to $A$. So I say that $A$, $C$ (summed) is also prime to $B$.
For since $DF$ is prime to each of $DE$, $EF$ then the (square) on
$DF$ is also prime to the (number created) from (multiplying) $DE$, $EF$
[Prop.~7.25]. But, the (sum of the squares) on $DE$, $EF$
plus twice the (number created) from (multiplying) $DE$, $EF$ is equal to
the (square) on $DF$ [Prop.~2.4]. And thus
the (sum of the squares) on $DE$, $EF$ plus twice  the (rectangle contained) by $DE$, $EF$ [is] prime to the (rectangle contained) by $DE$, $EF$.
By separation, the (sum of the squares) on $DE$, $EF$ plus  once the (rectangle contained) by $DE$, $EF$  is prime to the (rectangle contained)  by $DE$, $EF$.$^\dag$ Again, by separation,  the (sum of the squares) on $DE$, $EF$ is prime to the (rectangle contained)  by $DE$, $EF$. And the (square) on $DE$   is $A$,  and the (rectangle contained) by $DE$, $EF$ (is) $B$, and the (square) on $EF$ (is) $C$. Thus, $A$, $C$
summed is prime to $B$. (Which is) the very thing it was required to show.}
\end{Parallel}
{\footnotesize\noindent$^\dag$ Since if $\alpha\,\beta$ measures
$\alpha^2+\beta^2+2\,\alpha\,\beta$ then it also measures $\alpha^2+\beta^2+\alpha\,\beta$, and vice versa.}

%%%%%%
% Prop 9.16
%%%%%%
\pdfbookmark[1]{Proposition 9.16}{pdf9.16}
\begin{Parallel}{}{} 
\ParallelLText{
\begin{center}
{\large \ggn{16}.}
\end{center}\vspace*{-7pt}

{\greekfont Ἐὰν δύο ἀριθμοὶ πρῶτοι πρὸς ἀλλήλους ὦσιν, οὐκ ἔσται ὡς
ὁ πρῶτος πρὸς τὸν δεύτερον, οὕτως ὁ δεύτερος πρὸς ἄλλον τινά.}

\epsfysize=0.9in
\centerline{\epsffile{Book09/fig16g.eps}}

{\greekfont Δύο γὰρ ἀριθμοὶ οἱ Α, Β πρῶτοι πρὸς ἀλλήλους ἔστωσαν; λέγω,
ὅτι οὐκ ἔστιν ὡς ὁ Α πρὸς τὸν Β, οὕτως ὁ Β πρὸς ἄλλον
τινά.}

{\greekfont Εἰ γὰρ δυνατόν, ἔστω ὡς ὁ Α πρὸς τὸν Β, ὁ Β πρὸς τὸν Γ.
οἱ δὲ Α, Β πρῶτοι, οἱ δὲ πρῶτοι καὶ ἐλάθιστοι, οἱ δὲ ἐλάθιστοι
ἀριθμοὶ μετροῦσι τοὺς τὸν α ὐτὸν λόγον ἔθοντας ἰσάκις
ὅ τε ἡγούμενος τὸν ἡγούμενον καὶ ὁ ἑπόμενος τὸν ἑπόμενον;
μετρεῖ ἄρα ὁ Α τὸν Β ὡς ἡγούμενος ἡγούμενον.
μετρεῖ δὲ καὶ ἑαυτόν; ὁ Α ἄρα τοὺς Α, Β
μετρεῖ πρώτους ὄντας πρὸς ἀλλήλους; ὅπερ ἄτοπον. οὐκ ἄρα
ἔσται ὡς ὁ Α πρὸς τὸν Β, οὕτως ὁ Β πρὸς τὸν Γ; ὅπερ
ἔδει δεῖχαι.}}

\ParallelRText{
\begin{center}
{\large Proposition 16}
\end{center}

If two numbers are prime to one another then as the
first is to the second, so the second (will) not (be) to some other (number).

\epsfysize=0.9in
\centerline{\epsffile{Book09/fig16e.eps}}

For let the two numbers $A$ and $B$ be prime to one another. I say that
as $A$ is to $B$, so $B$ is not to some other (number).

For, if possible, let it be that as $A$ (is) to $B$, (so) $B$ (is) to $C$. And $A$ and $B$
(are) prime (to one another). And (numbers) prime (to one another are)
also the least (of those numbers having the same ratio as them) [Prop.~7.21]. And the least numbers measure those
(numbers) having the same ratio (as them) an equal number of times,
the leading (measuring) the leading, and the following the following [Prop.~7.20]. Thus, $A$ measures $B$, 
as the leading (measuring) the leading. And ($A$) also measures itself.
Thus, $A$ measures $A$ and $B$, which are prime to one another. The very
thing (is) absurd. Thus, as $A$  (is) to $B$, so $B$ cannot be to $C$.
(Which is) the very thing it was required to show.}
\end{Parallel}

%%%%%%
% Prop 9.17
%%%%%%
\pdfbookmark[1]{Proposition 9.17}{pdf9.17}
\begin{Parallel}{}{} 
\ParallelLText{
\begin{center}
{\large \ggn{17}.}
\end{center}\vspace*{-7pt}

{\greekfont Ἐὰν ὦσιν ὁσοιδηποτοῦν ἀριθμοὶ ἑχῆς ἀνάλογον, οἱ δὲ ἄκροι
α ὐτῶν πρῶτοι πρὸς ἀλλήλους ὦσιν, οὐκ ἔσται ὡς ὁ πρῶτος
πρὸς τὸν δεύτερον, οὕτως ὁ ἔσθατος πρὸς ἄλλον τινά.}

{\greekfont Ἔστωσαν ὁσοιδηποτοῦν ἀριθμοὶ ἑχῆς ἀνάλογον οἱ Α, Β, Γ, Δ,
οἱ δὲ ἄκροι α ὐτῶν οἱ Α, Δ πρῶτοι πρὸς ἀλλήλους ἔστωσαν;
λέγω, ὅτι οὐκ ἔστιν ὡς ὁ Α πρὸς τὸν Β, οὕτως ὁ Δ πρὸς ἄλλον
τινά.}\\

\epsfysize=1.3in
\centerline{\epsffile{Book09/fig17g.eps}}

{\greekfont Εἰ  γὰρ δυνατόν, ἔστω ὡς ὁ Α πρὸς τὸν Β, οὕτως
ὁ Δ πρὸς τὸν Ε; ἐναλλὰχ ἄρα ἐστὶν ὡς ὁ Α πρὸς τὸν Δ, ὁ Β πρὸς τὸν Ε. οἱ δὲ Α, Δ πρῶτοι, οἱ δὲ πρῶτοι καὶ ἐλάθιστοι, οἱ δὲ
ἐλάθιστοι ἀριθμοὶ μετροῦσι τοὺς τὸν α ὐτὸν λόγον ἔθοντας
ἰσάκις ὅ τε ἡγούμενος τὸν ἡγούμενον καὶ ὁ ἑπόμενος
τὸν ἑπόμενον. μετρεῖ ἄρα ὁ Α τὸν Β. καί ἐστιν ὡς ὁ Α
πρὸς τὸν Β, ὁ Β πρὸς τὸν Γ. καὶ ὁ Β ἄρα τὸν Γ μετρεῖ;
ὥστε καὶ ὁ Α τὸν Γ μετρεῖ. καὶ ἐπεί ἐστιν ὡς ὁ Β πρὸς
τὸν Γ, ὁ Γ πρὸς τὸν Δ, μετρεῖ δὲ ὁ Β τὸν Γ, μετρεῖ ἄρα καὶ
ὁ Γ τὸν Δ. ἀλλ ὁ Α τὸν Γ
ἐμέτρει; ὥστε
ὁ Α καὶ τὸν Δ μετρεῖ. μετρεῖ δὲ καὶ ἑαυτόν. ὁ Α ἄρα τοὺς
Α, Δ μετρεῖ πρώτους ὄντας πρὸς ἀλλήλους; ὅπερ ἐστὶν ἀδύνατον.
οὐκ ἄρα ἔσται ὡς ὁ Α πρὸς τὸν Β, οὕτως ὁ Δ πρὸς ἄλλον
τινά; ὅπερ ἔδει δεῖχαι.}}

\ParallelRText{
\begin{center}
{\large Proposition 17}
\end{center}

If   any multitude whatsoever of numbers is continuously proportional, and the outermost of them are prime to one another, then as the first (is) to the second, so the last will not be
to some other (number).

Let $A$, $B$, $C$, $D$ be any multitude whatsoever of continuously
proportional numbers. And let the outermost of them, $A$ and $D$,
be prime to one another. I say that as $A$ is to $B$, so $D$ (is) not to some
other (number).

\epsfysize=1.3in
\centerline{\epsffile{Book09/fig17e.eps}}

For, if possible, let it be that as $A$ (is) to $B$, so $D$ (is) to $E$. 
Thus, alternately, as $A$ is to $D$, (so) $B$ (is) to $E$ [Prop.~7.13]. And $A$ and $D$ are prime (to one another).  And (numbers) prime (to one another are)
also the least (of those numbers having the same ratio as them) [Prop.~7.21]. And the least numbers measure those
(numbers) having the same ratio (as them) an equal number of times,
the leading (measuring) the leading, and the following the following [Prop.~7.20]. Thus, $A$ measures $B$. And as $A$
is to $B$, (so) $B$ (is) to $C$.  Thus, $B$ also measures $C$.  And hence
$A$ measures $C$ [Def.~7.20]. And since as $B$ is to
$C$, (so) $C$ (is) to $D$, and $B$ measures $C$, $C$ thus also measures
$D$ [Def.~7.20].  But, $A$ was (found to be) measuring $C$. And 
hence $A$ also measures $D$. And ($A$) also measures itself.  Thus, $A$
measures $A$ and $D$, which are prime to one another. The very thing is impossible. Thus, as $A$ (is) to $B$, so $D$ cannot be to some other (number). (Which is) the very thing it was required to show.}
\end{Parallel}

%%%%%%
% Prop 9.18
%%%%%%
\pdfbookmark[1]{Proposition 9.18}{pdf9.18}
\begin{Parallel}{}{} 
\ParallelLText{
\begin{center}
{\large \ggn{18}.}
\end{center}\vspace*{-7pt}

{\greekfont Δύο ἀριθμῶν δοθέντων ἐπισκέψασθαι, εἰ δυνατόν ἐστιν α ὐτοῖς
τρίτον ἀνάλογον προσευρεῖν.}

\epsfysize=0.4in
\centerline{\epsffile{Book09/fig18g.eps}}

{\greekfont Ἔστωσαν οἱ δοθέντες δύο ἀριθμοὶ οἱ Α, Β, καὶ δέον ἔστω ἐπισκέψασθαι, εἰ δυνατόν ἐστιν α ὐτοῖς τρίτον ἀνάλογον προσευρεῖν.}

{\greekfont Οἱ  δὴ Α, Β ἤτοι πρῶτοι πρὸς ἀλλήλους εἰσὶν ἢ
οὔ. καὶ εἰ πρῶτοι πρὸς ἀλλήλους εἰσίν, δέδεικται, ὅτι
ἀδύνατόν ἐστιν α ὐτοῖς τρίτον ἀνάλογον προσευρεῖν.}

{\greekfont Ἀλλὰ δὴ μὴ ἔστωσαν οἱ Α, Β πρῶτοι πρὸς ἀλλήλους, καὶ ὁ Β ἑαυτον
πολλαπλασιάσας τὸν Γ ποιείτω. ὁ Α δὴ τὸν Γ ἤτοι μετρεῖ ἢ οὐ
μετρεῖ. μετρείτω πρότερον κατὰ τὸν Δ; ὁ Α ἄρα τὸν Δ πολλαπλασιάσας
τὸν Γ πεποίηκεν. ἀλλα μὴν καὶ ὁ Β ἑαυτὸν πολλαπλασιάσας τὸν Γ
πεποίηκεν; ὁ ἄρα ἐκ τῶν Α, Δ ἴσος ἐστὶ τῷ ἀπὸ τοῦ Β. ἔστιν
ἄρα ὡς ὁ Α πρὸς τὸν Β, ὁ Β πρὸς τὸν Δ; τοῖς Α, Β ἄρα τρίτος
ἀριθμὸς ἀνάλογον προσηύρηται ὁ Δ.}

{\greekfont Ἀλλὰ δὴ μὴ μετρείτω ὁ Α τὸν Γ; λέγω, ὅτι τοῖς Α, Β ἀδύνατόν ἐστι
τρίτον ἀνάλογον προσευρεῖν ἀριθμόν. εἰ γὰρ δυνατόν, προσηυρήσθω
ὁ Δ. ὁ ἄρα ἐκ τῶν Α, Δ ἴσος ἐστὶ τῷ ἀπὸ τοῦ Β. ὁ δὲ
ἀπὸ τοῦ Β ἐστιν ὁ Γ; ὁ ἄρα ἐκ τῶν
Α, Δ ἴσος ἐστὶ τῷ Γ. ὥστε ὁ Α τὸν Δ πολλαπλασιάσας τὸν Γ 
πεποίηκεν; ὁ Α ἄρα τὸν Γ μετρεῖ κατὰ τὸν Δ. ἀλλα μὴν ὑπόκειται
καὶ μὴ μετρῶν; ὅπερ ἄτοπον. οὐκ ἄρα δυνατόν ἐστι
τοῖς Α, Β τρίτον ἀνάλογον προσευρεῖν ἀριθμὸν, ὅταν ὁ Α τὸν Γ μὴ
μετρῇ; ὅπερ ἔδει δεῖχαι.}}

\ParallelRText{
\begin{center}
{\large Proposition 18}
\end{center}

For  two given numbers, to investigate whether it
is possible to find a third (number) proportional to them.

\epsfysize=0.4in
\centerline{\epsffile{Book09/fig18e.eps}}

Let $A$ and $B$ be the two given numbers. And let it be required to investigate whether it is possible to find a third (number) proportional
to them.

So $A$ and $B$ are either prime to one another, or not. 
And if they are prime to one another then it has (already) been show that it is impossible
to find a third (number) proportional to them [Prop.~9.16].

And so let $A$ and $B$ not be prime to one another. And let $B$ make $C$
(by) multiplying itself. So $A$ either measures, or does not measure, $C$. 
Let it first of all measure ($C$) according to $D$. Thus, $A$ has made
$C$ (by) multiplying $D$. But, in fact, $B$ has also made
$C$ (by) multiplying itself. Thus, the (number created) from (multiplying)
$A$, $D$ is equal to the (square) on $B$. Thus, as $A$ is to $B$, (so)
$B$ (is) to $D$ [Prop.~7.19]. Thus, a third number
has been found proportional to $A$, $B$, (namely) $D$. 

And so let $A$ not measure $C$. I say that it is impossible to find a third
number proportional to $A$, $B$. For, if possible, let it be found,
(and let it be) $D$. Thus, the (number created) from (multiplying) $A$, $D$
is equal to the (square) on $B$ [Prop.~7.19]. 
And the (square) on $B$ is $C$. Thus, the (number created) from (multiplying) $A$, $D$
is equal to $C$. Hence, $A$ has made $C$ (by) multiplying $D$. 
Thus, $A$ measures $C$ according to $D$. But ($A$) was, in fact, also assumed
(to be) not measuring ($C$). The very thing (is) absurd. Thus, it is not
possible to find a third number proportional to $A$, $B$ when $A$ does not
measure $C$. (Which is) the very thing it was required to show.}
\end{Parallel}

%%%%%%
% Prop 9.19
%%%%%%
\pdfbookmark[1]{Proposition 9.19}{pdf9.19}
\begin{Parallel}{}{} 
\ParallelLText{
\begin{center}
{\large \ggn{19}.}
\end{center}\vspace*{-7pt}

{\greekfont Τριῶν ἀριθμῶν δοθέντων ἐπισκέψασθαι, πότε δυνατόν ἐστιν α ὐτοῖς
τέταρτον ἀνάλογον προσευρεῖν.}

\epsfysize=1.3in
\centerline{\epsffile{Book09/fig19g.eps}}

{\greekfont Ἔστωσαν οἱ δοθέντες τρεῖς ἀριθμοὶ οἱ Α, Β, Γ, καὶ δέον ἔστω
επισκέψασθαι, πότε δυνατόν ἐστιν α ὐτοῖς τέταρτον
ἀνάλογον προσευρεῖν.}

{\greekfont Ἤτοι οὖν οὔκ εἰσιν ἑχῆς ἀνάλογον, καὶ οἱ ἄκροι α ὐτῶν πρῶτοι
πρὸς ἀλλήλους εἰσίν, ἢ ἑχῆς εἰσιν ἀνάλογον, καὶ οἱ
ἄκροι α ὐτῶν οὔκ εἰσι πρῶτοι πρὸς ἀλλήλους, ἢ οὕτε
ἑχῆς εἰσιν ἀνάλογον, οὔτε οἱ ἄκροι α ὐτῶν πρῶτοι
πρὸς ἀλλήλους εἰσίν, ἢ καὶ ἑχῆς εἰσιν ἀνάλογον, καὶ
οἱ ἄκροι α ὐτῶν πρῶτοι πρὸς ἀλλήλους εἰσίν.}

{\greekfont Εἰ μὲν οὖν οἱ Α, Β, Γ ἑχῆς εἰσιν ἀνάλογον, καὶ οἱ ἄκροι
α ὐτῶν οἱ Α, Γ πρῶτοι πρὸς ἀλλήλους εἰσίν, δέδεικται, ὅτι
ἀδύνατόν ἐστιν α ὐτοῖς τέταρτον ἀνάλογον προσευρεῖν ἀριθμόν. μὴ
ἔστωσαν δὴ οἱ Α, Β, Γ ἑχῆς ἀνάλογον τῶν ἀκρῶν πάλιν
ὄντων πρώτων πρὸς  ἀλλήλους. λέγω, ὅτι καὶ οὕτως ἀδύνατόν ἐστιν
α ὐτοῖς τέταρτον ἀνάλογον προσευρεῖν. εἰ γὰρ δυνατόν, προσευρήσθω
ὁ Δ, ὥστε εἶναι ὡς τὸν Α πρὸς τὸν Β, τὸν Γ πρὸς τὸν Δ, καὶ
γεγονέτω ὡς ὁ Β πρὸς τὸν Γ, ὁ Δ πρὸς τὸν Ε. καὶ ἐπεί ἐστιν ὡς μὲν ὁ Α πρὸς τὸν Β, ὁ Γ πρὸς τὸν Δ, ὡς  δὲ ὁ Β πρὸς τὸν 
Γ, ὁ Δ πρὸς τὸν Ε, δι ἴσου ἄρα ὡς ὁ Α πρὸς τὸν Γ, ὁ Γ πρὸς τὸν Ε. οἱ δὲ Α, Γ πρῶτοι, οἱ δὲ
πρῶτοι καὶ ἐλάθιστοι, οἱ δὲ ἐλάθιστοι μετροῦσι τοὺς τὸν α ὐτὸν
λόγον ἔθοντας ὅ τε ἡγούμενος τὸν ἡγούμενον
καὶ ὁ ἑπόμενος τὸν ἑπόμενον. μετρεῖ ἄρα ὁ Α τὸν Γ ὡς
ἡγούμενος ἡγούμενον. μετρεῖ δὲ καὶ ἑαυτόν;
ὁ Α ἄρα τοὺς Α, Γ μετρεῖ πρώτους ὄντας πρὸς ἀλλήλους; ὅπερ
ἐστὶν ἀδύνατον. οὐκ ἄρα τοῖς Α, Β, Γ δυνατόν ἐστι τέταρτον
ἀνάλογον προσευρεῖν.}

{\greekfont Ἀλλά δὴ πάλιν ἔστωσαν οἱ Α, Β, Γ ἑχῆς ἀνάλογον, οἱ δὲ Α, Γ μὴ
ἔστωσαν πρῶτοι πρὸς ἀλλήλους. λέγω, ὅτι δυνατόν ἐστιν α ὐτοῖς τέταρτον ἀνάλογον προσευρεῖν. ὁ γὰρ Β τὸν Γ πολλαπλασιάσας
τὸν Δ ποιείτω; ὁ Α ἄρα τὸν Δ ἤτοι μετρεῖ ἢ οὐ μετρεῖ. μετρείτω
α ὐτὸν πρότερον κατὰ τὸν Ε; ὁ Α ἄρα τὸν Ε πολλαπλασιάσας τὸν Δ
πεποίηκεν. ἀλλὰ μὴν καὶ ὁ Β τὸν Γ  πολλαπλασιάσας τὸν Δ πεποίηκεν;
ὁ ἄρα ἐκ τῶν Α, Ε ἴσος ἐστὶ τῷ ἐκ τῶν Β, Γ. ἀνάλογον
ἄρα [ἐστὶν] ὡς ὁ Α πρὸς τὸν Β, ὁ Γ πρὸς τὸν Ε; τοὶς Α, Β, Γ
ἄρα τέταρτος ἀνάλογον προσηύρηται ὁ Ε.}

{\greekfont Ἀλλὰ δὴ μὴ μετρείτω ὁ Α τὸν Δ; λέγω, ὅτι ἀδύνατόν ἐστι τοῖς
Α, Β, Γ τέταρτον ἀνάλογον προσευρεῖν ἀριθμόν. εἰ γὰρ δυνατόν,
προσευρήσθω ὁ Ε; ὁ ἄρα ἐκ τῶν Α, Ε ἴσος ἐστὶ τῷ ἐκ τῶν Β, Γ.
ἀλλὰ ὁ ἑκ τῶν Β, Γ ἐστιν ὁ Δ; καὶ ὁ ἐκ τῶν Α, Ε ἄρα ἴσος ἐστὶ τῷ Δ. ὁ Α ἄρα τὸν Ε πολλαπλασιάσας τὸν Δ πεποίηκεν; ὁ Α
ἄρα τὸν Δ μετρεῖ κατὰ τὸν  Ε; ὥστε μετρεῖ ὁ Α τὸν Δ. ἀλλὰ
καὶ οὐ μετρεῖ; ὅπερ ἄτοπον. οὐκ ἄρα δυνάτον ἐστι τοῖς
Α, Β, Γ τέταρτον ἀνάλογον προσευρεῖν ἀριθμόν, ὅταν ὁ Α τὸν
Δ μὴ μετρῇ. ἀλλὰ δὴ οἱ Α, Β, Γ μήτε ἑχῆς ἔστωσαν ἀνάλογον
μήτε οἱ ἄκροι πρῶτοι πρὸς ἀλλήλους. καὶ ὁ Β τὸν Γ πολλαπλασιάσας
τὸν Δ ποιείτω. ὁμοίως δὴ δειθθήσεται, ὅτι εἰ μὲν μετρεῖ ὁ Α τὸν Δ,
δυνατόν ἐστιν α ὐτοῖς ἀνάλογον προσευρεῖν, εἰ δὲ οὐ μετρεῖ,
ἀδύνατον; ὅπερ ἔδει δεῖχαι.}}

\ParallelRText{
\begin{center}
{\large Proposition 19}$^\dag$
\end{center}

For three given numbers, to investigate when it is
possible to find a fourth (number) proportional to them.

\epsfysize=1.3in
\centerline{\epsffile{Book09/fig19e.eps}}

Let $A$, $B$, $C$ be the three given numbers. And let it be required to
investigate when it is possible to find a fourth (number) proportional to them.

In fact, ($A$, $B$, $C$) are either not continuously proportional and the outermost
of them are prime to one another, or are continuously proportional
and the outermost of them are not prime to one another, or are neither
continuously proportional nor are the outermost of them prime to one another, or are continuously proportional and the outermost of them are
prime to one another.

In fact, if $A$, $B$, $C$ are continuously proportional, and the outermost
of them, $A$ and $C$, are prime to one another, (then) it has (already) been shown that it is impossible to find a fourth number proportional to them
[Prop.~9.17]. So let $A$, $B$, $C$
not be continuously proportional, (with) the outermost of them again being
prime to one another. I say that, in this case, it is also impossible to find a fourth (number)
proportional to them. For, if possible, let it be found,
(and let it be) $D$. Hence, it will be that as $A$ (is) to $B$, (so)
$C$ (is) to $D$. And let it be contrived that as $B$ (is) to $C$, (so) $D$ (is) to $E$. And since as $A$ is to $B$, (so) $C$ (is) to $D$, and as $B$ (is) to $C$, (so)
$D$ (is) to $E$, thus, via equality, as $A$ (is) to $C$, (so) $C$ (is) to $E$
[Prop.~7.14]. And $A$ and $C$ (are) prime (to one another). And (numbers) prime (to one another are) also the least (numbers
having the same ratio as them) [Prop.~7.21]. 
And the least (numbers) measure those numbers having the same ratio
as them (the same number of times), the leading (measuring) the leading,
and the following the following [Prop.~7.20]. 
Thus, $A$ measures $C$, (as) the leading (measuring) the leading.
And it also measures itself. Thus, $A$ measures $A$ and $C$, which are
prime to one another. The very thing is impossible. Thus, it is not possible to find a fourth (number) proportional to $A$, $B$, $C$.

And so let $A$, $B$, $C$  again be continuously proportional, and let $A$ and $C$ not be prime to one another. I say that it is possible to find
a fourth (number) proportional to them. For let $B$ make $D$ (by) multiplying $C$. Thus, $A$ either measures or does not measure $D$. 
Let it, first of all, measure ($D$) according to $E$. Thus, $A$ has made $D$ (by) multiplying $E$. But, in fact, $B$ has also made $D$ (by) multiplying $C$. Thus, the (number created) from (multiplying) $A$, $E$
is equal to the (number created) from (multiplying) $B$, $C$. Thus, 
proportionally, as $A$ [is] to $B$, (so) $C$ (is) to $E$ [Prop.~7.19]. Thus, a fourth (number) proportional
to $A$, $B$, $C$ has been found, (namely) $E$.

And so let $A$ not measure $D$. I say that it is impossible to find a fourth
number proportional to $A$, $B$, $C$. For, if possible, let it 
be found, (and let it be) $E$. Thus, the (number created) from
(multiplying) $A$, $E$ is equal to the (number created) from (multiplying)
$B$, $C$. But, the (number created) from (multiplying) $B$, $C$ is $D$.
And thus the (number created) from (multiplying) $A$, $E$ is equal to $D$.
Thus, $A$ has made $D$ (by) multiplying $E$. Thus, $A$ measures
$D$ according to $E$. Hence, $A$ measures $D$. But, it also does not
measure ($D$). The very thing (is) absurd. Thus, it is not possible
to find a fourth number proportional to $A$, $B$, $C$ when $A$ does not
measure $D$. And so (let) $A$, $B$, $C$ (be) neither continuously
proportional, nor (let) the outermost of them (be) prime to one another.
And let $B$ make $D$ (by) multiplying $C$. So, similarly, it can be show that
if $A$ measures $D$ then it is possible to find a fourth (number)
proportional to ($A$, $B$, $C$), and impossible if ($A$) does not measure ($D$).
(Which is) the very thing it was required to show.}
\end{Parallel}
{\footnotesize\noindent$^\dag$ The proof of this proposition is incorrect.
There are, in fact, only two cases. Either $A$, $B$, $C$ are continuously
proportional,  with $A$ and $C$ prime to one another, or not. In the first case, 
it is impossible to find a fourth proportional number. In the second case,
it is possible to find a fourth proportional number provided that $A$
measures $B$ times $C$. Of the four cases considered by Euclid,
the proof given in the second case is incorrect, since it only demonstrates that
if $A:B::C:D$ then a number $E$ cannot be found such that $B:C::D:E$. The proofs given in the other three cases are correct.}

%%%%%%
% Prop 9.20
%%%%%%
\pdfbookmark[1]{Proposition 9.20}{pdf9.20}
\begin{Parallel}{}{} 
\ParallelLText{
\begin{center}
{\large\ggn{20}.}
\end{center}\vspace*{-7pt}

{\greekfont Οἱ πρῶτοι ἀριθμοὶ πλείους εἰσὶ παντὸς τοῦ προτεθέντος πλήθους
πρώτων ἀριθμῶν.}

\epsfysize=1.3in
\centerline{\epsffile{Book09/fig20g.eps}}

{\greekfont Ἔστωσαν οἱ προτεθέντες πρῶτοι ἀριθμοὶ οἱ Α, Β, Γ; λέγω,
ὅτι τῶν Α, Β, Γ πλείους εἰσὶ πρῶτοι ἀριθμοί.}

{\greekfont Εἰλήφθω γὰρ ὁ ὑπὸ τῶν Α, Β, Γ ἐλάθιστος μετρούμενος καὶ ἔστω
ΔΕ, καὶ προσκείσθω τῷ ΔΕ μονὰς ἡ ΔΖ. ὁ δὴ ΕΖ ἤτοι πρῶτός
ἐστιν ἢ οὔ. ἔστω πρότερον πρῶτος; εὐρημένοι ἄρα εἰσὶ πρῶτοι
ἀριθμοὶ οἱ Α, Β, Γ, ΕΖ πλείους τῶν Α, Β, Γ.}

{\greekfont Ἀλλὰ δὴ μὴ ἔστω ὁ ΕΖ πρῶτος; ὑπὸ πρώτου
ἄρα τινὸς ἀριθμοῦ μετρεῖται. μετρείσθω ὑπὸ πρώτου τοῦ Η;
λέγω, ὅτι ὁ Η οὐδενὶ τῶν Α, Β, Γ ἐστιν ὁ α ὐτός. εἰ γὰρ
δυνατόν, ἔστω. οἱ δὲ Α, Β, Γ τὸν ΔΕ μετροῦσιν; καὶ ὁ Η ἄρα
τὸν ΔΕ μετρήσει. μετρεῖ δὲ καὶ τὸν ΕΖ; καὶ λοιπὴν τὴν ΔΖ μονάδα
μετρήσει ὁ Η ἀριθμὸς ὤν; ὄπερ ἄτοπον. οὐκ ἄρα ὁ
Η ἑνὶ τῶν Α, Β, Γ ἐστιν ὁ α ὐτός. καὶ ὑπόκειται πρῶτος.
εὑρημένοι ἄρα εἰσὶ πρῶτοι ἀριθμοὶ πλείους
τοῦ προτεθέντος πλήθους τῶν Α, Β, Γ οἱ Α, Β, Γ, Η;
ὅπερ ἔδει δεῖχαι.}}

\ParallelRText{
\begin{center}
{\large Proposition 20}
\end{center}

The (set of all) prime numbers is more numerous than any assigned
multitude of prime numbers.

\epsfysize=1.3in
\centerline{\epsffile{Book09/fig20e.eps}}

Let $A$, $B$, $C$ be the assigned prime numbers. I say that the (set of all)
primes numbers is more numerous than $A$, $B$, $C$.

For let the least number measured by  $A$, $B$, $C$ be taken, and let it be $DE$ [Prop.~7.36]. And let the unit $DF$ be added to $DE$. So $EF$ is either prime, or not. Let it, first of all,
be prime. Thus, the (set of) prime numbers $A$, $B$, $C$, $EF$, (which is) more numerous than $A$, $B$, $C$, has been found.

And so let $EF$ not be prime. Thus, it is measured by some prime number [Prop.~7.31]. Let it be measured by the prime
(number) $G$. I say that $G$ is not the same as any of $A$, $B$, $C$. 
For, if possible, let it be (the same). And $A$, $B$, $C$ (all) measure
$DE$. Thus, $G$ will also measure $DE$. And it also measures $EF$. (So)
$G$ will also measure the remainder, unit $DF$, (despite) being a number
[Prop.~7.28]. The very thing (is) absurd. Thus, $G$
is not the same as one of $A$, $B$, $C$. And it was assumed (to be) prime.
Thus, the (set of) prime numbers $A$, $B$, $C$, $G$, (which is) more numerous than the assigned multitude (of prime numbers), $A$, $B$, $C$, has been found. (Which is) the very thing it
was required to show.}
\end{Parallel}

%%%%%%
% Prop 9.21
%%%%%%
\pdfbookmark[1]{Proposition 9.21}{pdf9.21}
\begin{Parallel}{}{} 
\ParallelLText{
\begin{center}
{\large \ggn{21}.}
\end{center}\vspace*{-7pt}

{\greekfont Ἐὰν ἄρτιοι ἀριθμοὶ ὁποσοιοῦν συντεθῶσιν, ὁ ὅλος ἄρτιός
ἐστιν.}

\epsfysize=0.3in
\centerline{\epsffile{Book09/fig21g.eps}}

{\greekfont Συγκείσθωσαν γὰρ ἄρτιοι ἀριθμοὶ ὁποσοιοῦν οἱ ΑΒ, ΒΓ, ΓΔ, ΔΕ;
λέγω, ὅτι ὅλος ὁ ΑΕ ἄρτιός ἐστιν.}

{\greekfont Ἐπεὶ γὰρ ἕκαστος τῶν ΑΒ, ΒΓ, ΓΔ, ΔΕ ἄρτιός ἐστιν, ἔχει
μέρος ἥμισυ; ὥστε καὶ ὅλος ὁ ΑΕ ἔχει μέρος ἥμισυ.
ἄρτιος δὲ ἀριθμός ἐστιν ὁ δίθα διαιρούμενος;
ἄρτιος ἄρα ἐστὶν ὁ ΑΕ; ὅπερ ἔδει δεῖχαι.}}

\ParallelRText{
\begin{center}
{\large Proposition 21}
\end{center}

If any multitude whatsoever of even numbers
is added together then the whole is even.

\epsfysize=0.3in
\centerline{\epsffile{Book09/fig21e.eps}}

For let any multitude  whatsoever of even numbers, $AB$, $BC$, $CD$, $DE$, lie together. I say that the whole, $AE$, is even.

For since everyone  of $AB$, $BC$, $CD$, $DE$ is even, it has a half part
[Def.~7.6]. And hence the whole $AE$ has
a half part. And an even number is one (which can be) divided in half [Def.~7.6]. Thus, $AE$ is even. (Which is) the very thing it was required to show.}
\end{Parallel}

%%%%%%
% Prop 9.22
%%%%%%
\pdfbookmark[1]{Proposition 9.22}{pdf9.22}
\begin{Parallel}{}{} 
\ParallelLText{
\begin{center}
{\large \ggn{22}.}
\end{center}\vspace*{-7pt}

{\greekfont Ἐὰν περισσοὶ ἀριθμοὶ ὁποσοιοῦν συντεθῶσιν, τὸ δὲ πλῆθος
α ὐτῶν ἄρτιον ᾖ, ὁ ὅλος ἄρτιος ἔσται.}\\

\epsfysize=0.3in
\centerline{\epsffile{Book09/fig21g.eps}}

{\greekfont Συγκείσθωσαν γὰρ περισσοὶ ἀριθμοὶ ὁσοιδηποτοῦν ἄρτιοι
τὸ πλῆθος οἱ ΑΒ, ΒΓ, ΓΔ, ΔΕ; λέγω, ὅτι ὅλος ὁ ΑΕ ἄρτιός
ἐστιν.}

{\greekfont Ἐπεὶ γὰρ ἕκαστος τῶν ΑΒ, ΒΓ, ΓΔ, ΔΕ περιττός ἐστιν, ἀφαιρεθείσης
μονάδος ἀφ ἑκάστου ἕκαστος τῶν λοιπῶν ἄρτιος ἔσται; ὥστε
καὶ ὁ συγκείμενος ἐχ α ὐτῶν ἄρτιος ἔσται. ἔστι δὲ καὶ τὸ
πλῆθος τῶν μονάδων ἄρτιον. καὶ ὅλος ἄρα ὁ ΑΕ ἄρτιός
ἐστιν; ὅπερ ἔδει δεῖχαι.}}

\ParallelRText{
\begin{center}
{\large Proposition 22}
\end{center}

If any multitude whatsoever of odd numbers
is added together, and the multitude of them is even,  then the whole will be even.

\epsfysize=0.3in
\centerline{\epsffile{Book09/fig21e.eps}}

For let any even multitude whatsoever of odd numbers, 
$AB$, $BC$, $CD$, $DE$, lie together. I say that the whole, $AE$, is even.

For since everyone of $AB$, $BC$, $CD$, $DE$ is odd then,  a unit being subtracted from
each, everyone of the remainders will be (made) even [Def.~7.7]. And hence the sum of them will be even [Prop.~9.21]. And the multitude  of the units is even.
Thus, the whole $AE$ is also even [Prop.~9.21]. (Which is) the very thing it was required to show.}
\end{Parallel}

%%%%%%
% Prop 9.23
%%%%%%
\pdfbookmark[1]{Proposition 9.23}{pdf9.23}
\begin{Parallel}{}{} 
\ParallelLText{
\begin{center}
{\large\ggn{23}.}
\end{center}\vspace*{-7pt}

{\greekfont Ἐὰν περισσοὶ ἀριθμοὶ ὁποσοιοῦν συντεθῶσιν, τὸ
δὲ πλῆθος α ὐτῶν περισσὸν ᾖ, καὶ ὁ ὅλος περισσὸς ἔσται.}\\

\epsfysize=0.3in
\centerline{\epsffile{Book09/fig23g.eps}}

{\greekfont Συγκείσθωσαν γὰρ ὁποσοιοῦν περισσοὶ ἀριθμοί, ὧν τὸ πλῆθος
περισσὸν ἔστω, οἱ ΑΒ, ΒΓ, ΓΔ; λέγω, ὅτι καὶ ὅλος ὁ ΑΔ περισσός
ἐστιν.}

{\greekfont Ἀφῃρήσθω ἀπὸ τοῦ ΓΔ μονὰς ἡ ΔΕ; λοιπὸς ἄρα ὁ ΓΕ ἄρτιός
ἐστιν. ἔστι δὲ καὶ ὁ ΓΑ ἄρτιος; καὶ ὅλος ἄρα ὁ ΑΕ ἄρτιός 
ἐστιν. καί ἐστι μονὰς ἡ ΔΕ. περισσὸς ἄρα ἐστὶν ὁ ΑΔ; ὅπερ
ἔδει δεῖχαι.}}

\ParallelRText{
\begin{center}
{\large Proposition 23}
\end{center}

If  any multitude whatsoever of odd numbers
is added together, and the multitude of them is odd,  then the whole will also be odd.

\epsfysize=0.3in
\centerline{\epsffile{Book09/fig23e.eps}}

For let any  multitude whatsoever of odd numbers, $AB$, $BC$, $CD$,  lie together, and let the multitude of them be odd. I say that the whole, $AD$, is also odd.
 
For let the unit $DE$ be subtracted from $CD$. The remainder, $CE$,
is thus even [Def.~7.7]. And $CA$ is also even
[Prop.~9.22]. Thus, the whole, $AE$, is also even
[Prop.~9.21]. And $DE$ is a unit. Thus,
$AD$ is odd [Def.~7.7]. (Which is) the very thing it
was required to show.}
\end{Parallel}

%%%%%%
% Prop 9.24
%%%%%%
\pdfbookmark[1]{Proposition 9.24}{pdf9.24}
\begin{Parallel}{}{} 
\ParallelLText{
\begin{center}
{\large \ggn{24}.}
\end{center}\vspace*{-7pt}

{\greekfont Ἐὰν ἀπὸ ἀρτίου ἀριθμοῦ ἄρτιος ἀφαιρεθῇ, ὁ λοιπὸς
ἄρτιος ἔσται.}

\epsfysize=0.3in
\centerline{\epsffile{Book09/fig24g.eps}}

{\greekfont Ἀπὸ γὰρ ἀρτίου τοῦ ΑΒ ἄρτιος ἀφῃρήσθω ὁ ΒΓ; λέγω, ὅτι
ὁ λοιπὸς ὁ ΓΑ ἄρτιός ἐστιν.}

{\greekfont Ἐπεὶ γὰρ ὁ ΑΒ ἄρτιός ἐστιν, ἔχει μέρος ἥμισυ.
διὰ τὰ α ὐτὰ δὴ καὶ ὁ ΒΓ ἔχει μέρος ἥμισυ; ὥστε
καὶ λοιπὸς [ὁ ΓΑ ἔχει μέρος ἥμισυ] ἄρτιος [ἄρα]
ἐστὶν ὁ ΑΓ; ὅπερ ἔδει δεῖχαι.}}

\ParallelRText{
\begin{center}
{\large Proposition 24}
\end{center}

If  an even (number) is subtracted from an(other) even
number then the remainder will be even.

\epsfysize=0.3in
\centerline{\epsffile{Book09/fig24e.eps}}

For let the even (number) $BC$ be subtracted from the even number
$AB$. I say that the remainder, $CA$, is even.

For since $AB$ is even, it has a half part [Def.~7.6]. 
So, for the same (reasons), $BC$ also has a half part. And hence the remainder [$CA$ has a half part]. [Thus,] $AC$ is even. (Which is) the very thing it was required to show.}
\end{Parallel}

%%%%%%
% Prop 9.25
%%%%%%
\pdfbookmark[1]{Proposition 9.25}{pdf9.25}
\begin{Parallel}{}{} 
\ParallelLText{
\begin{center}
{\large \ggn{25}.}
\end{center}\vspace*{-7pt}

{\greekfont Ἐὰν ἀπὸ ἀρτίου ἀριθμοῦ περισσὸς ἀφαιρεθῇ, ὁ λοιπὸς
περισσὸς ἔσται.}

\epsfysize=0.3in
\centerline{\epsffile{Book09/fig25g.eps}}

{\greekfont Ἀπὸ γὰρ ἀρτίου τοῦ ΑΒ περισσὸς ἀφῃρήσθω ὁ ΒΓ; λέγω, ὅτι
ὁ λοιπὸς ὁ ΓΑ περισσός ἐστιν.}

{\greekfont Ἀφῃρήσθω γὰρ ἀπὸ τοῦ ΒΓ μονὰς ἡ ΓΔ; ὁ ΔΒ ἄρα ἄρτιός
ἐστιν. ἔστι δὲ καὶ ὁ ΑΒ ἄρτιος; καὶ λοιπὸς ἄρα ὁ ΑΔ ἄρτιός
ἐστιν. καί ἐστι μονὰς ἡ ΓΔ; ὁ ΓΑ ἄρα περισσός ἐστιν; ὅπερ
ἔδει δεῖχαι.}}

\ParallelRText{
\begin{center}
{\large Proposition 25}
\end{center}

If an odd (number) is subtracted from an even
number then the remainder will be odd.

\epsfysize=0.3in
\centerline{\epsffile{Book09/fig25e.eps}}

For let the odd (number) $BC$ be subtracted from the even number
$AB$. I say that the remainder $CA$ is odd.

For let the unit $CD$ be subtracted from $BC$. $DB$ is thus even
[Def.~7.7]. And $AB$ is also even. And thus the remainder $AD$ is even [Prop.~9.24]. And $CD$ is a unit. 
Thus, $CA$ is odd [Def.~7.7]. (Which is) the very thing it was required to show.}
\end{Parallel}

%%%%%%
% Prop 9.26
%%%%%%
\pdfbookmark[1]{Proposition 9.26}{pdf9.26}
\begin{Parallel}{}{} 
\ParallelLText{
\begin{center}
{\large \ggn{26}.}
\end{center}\vspace*{-7pt}

{\greekfont Ἐὰν ἀπὸ  περισσοῦ ἀριθμοῦ περισσὸς ἀφαιρεθῇ, ὁ λοιπὸς ἄρτιος
ἔσται.}

\epsfysize=0.3in
\centerline{\epsffile{Book09/fig26g.eps}}

{\greekfont Ἀπὸ γὰρ περισσοῦ τοῦ ΑΒ περισσὸς ἀφῃρήσθω ὁ ΒΓ; λέγω, ὅτι
ὁ λοιπὸς ὁ ΓΑ ἄρτιός ἐστιν.}

{\greekfont Ἐπεὶ γὰρ ὁ ΑΒ περισσός ἐστιν, ἀφῃρήσθω μονὰς ἡ ΒΔ; λοιπὸς
ἄρα ὁ ΑΔ ἄρτιός ἐστιν. διὰ τὰ α ὐτὰ δὴ καὶ ὁ ΓΔ ἄρτιός
ἐστιν; ὥστε καὶ λοιπὸς ὁ ΓΑ ἄρτιός ἐστιν; ὅπερ ἔδει
δεῖχαι.}}

\ParallelRText{
\begin{center}
{\large Proposition 26}
\end{center}

If an odd (number) is subtracted from an odd
number then the remainder will be even.

\epsfysize=0.3in
\centerline{\epsffile{Book09/fig26e.eps}}

For let the odd (number) $BC$ be subtracted from the odd
(number) $AB$. I say that the remainder, $CA$, is even.

For since $AB$ is odd, let the unit $BD$ be subtracted (from it). 
Thus, the remainder $AD$ is even [Def.~7.7]. So, for the same (reasons), 
$CD$ is also even. And hence the remainder $CA$ is even [Prop.~9.24]. (Which is) the very thing it was required to show.}
\end{Parallel}

%%%%%%
% Prop 9.27
%%%%%%
\pdfbookmark[1]{Proposition 9.27}{pdf9.27}
\begin{Parallel}{}{} 
\ParallelLText{
\begin{center}
{\large \ggn{27}.}
\end{center}\vspace*{-7pt}

{\greekfont Ἐὰν ἀπὸ περισσοῦ ἀριθμοῦ ἄρτιος ἀφαιρεθῇ, ὁ λοιπὸς περισσὸς 
ἔσται.}

\epsfysize=0.3in
\centerline{\epsffile{Book09/fig27g.eps}}

{\greekfont Ἀπὸ γὰρ περισσοῦ τοῦ ΑΒ ἄρτιος ἀφῃρήσθω ὁ ΒΓ; λέγω, ὅτι
ὁ λοιπὸς ὁ ΓΑ περισσός ἐστιν.}

{\greekfont Ἀφῃρήσθω [γὰρ] μονὰς ἡ ΑΔ; ὁ ΔΒ ἄρα ἄρτιός ἐστιν. ἔστι
δὲ καὶ ὁ ΒΓ ἄρτιος; καὶ λοιπὸς ἄρα ὁ ΓΔ ἄρτιός ἐστιν.
περισσὸς ἄρα ὁ ΓΑ; ὅπερ ἔδει δεῖχαι.}}

\ParallelRText{
\begin{center}
{\large Proposition 27}
\end{center}

If an even (number) is subtracted from an odd
number then the remainder will be odd.

\epsfysize=0.3in
\centerline{\epsffile{Book09/fig27e.eps}}

For let the even (number) $BC$ be subtracted from the odd (number)
$AB$. I say that the remainder, $CA$, is odd.

$[$For$]$ let the unit $AD$ be subtracted (from $AB$). $DB$ is thus
even [Def.~7.7]. And $BC$ is also even. Thus, the
remainder $CD$ is also even [Prop.~9.24]. 
$CA$ (is) thus odd [Def.~7.7]. (Which is) the
very thing it was required to show.}
\end{Parallel}

%%%%%%
% Prop 9.28
%%%%%%
\pdfbookmark[1]{Proposition 9.28}{pdf9.28}
\begin{Parallel}{}{} 
\ParallelLText{
\begin{center}
{\large \ggn{28}.}
\end{center}\vspace*{-7pt}

{\greekfont Ἐὰν περισσὸς ἀριθμὸς ἄρτιον πολλαπλασιάσας ποιῇ τινα, ὁ γενόμενος
ἄρτιος ἔσται.}

\epsfysize=0.8in
\centerline{\epsffile{Book09/fig28g.eps}}

{\greekfont Περισσὸς γὰρ ἀριθμὸς ὁ Α ἄρτιον τὸν Β πολλαπλασιάσας τὸν Γ ποιείτω;
λέγω, ὅτι ὁ Γ ἄρτιός ἐστιν.}

{\greekfont Ἐπεὶ γὰρ ὁ Α τὸν Β πολλαπλασιάσας τὸν Γ πεποίηκεν, ὁ Γ ἄρα σύγκειται ἐκ τοσούτων ἴσων τῷ Β, 
ὅσαι εἰσὶν ἐν τῷ Α μονάδες. καί ἐστιν ὁ Β ἄρτιος;
ὁ Γ ἄρα σύγκειται ἐχ ἀρτίων. ἐὰν δὲ ἄρτιοι ἀριθμοὶ
ὁποσοιοῦν συντεθῶσιν, ὁ ὅλος ἄρτιός ἐστιν. ἄρτιος
ἄρα ἐστὶν ὁ Γ; ὅπερ ἔδει δεῖχαι.}}

\ParallelRText{
\begin{center}
{\large Proposition 28}
\end{center}

If an odd number makes some (number by)
multiplying an even (number) then the created (number) will be even.

\epsfysize=0.8in
\centerline{\epsffile{Book09/fig28e.eps}}

For let the odd number $A$ make $C$ (by) multiplying the even (number)
$B$. I say that $C$ is even.

For since $A$ has made $C$ (by) multiplying $B$, $C$ is thus composed  out of so many (magnitudes) equal to $B$, as many as (there) are  units in $A$ [Def.~7.15]. And $B$ is even. Thus, $C$ is composed
out of even (numbers). And if any multitude whatsoever of even numbers
is added together then the whole is even [Prop.~9.21]. 
Thus, $C$ is even. (Which is) the very thing it was required to show.}
\end{Parallel}

%%%%%%
% Prop 9.29
%%%%%%
\pdfbookmark[1]{Proposition 9.29}{pdf9.29}
\begin{Parallel}{}{} 
\ParallelLText{
\begin{center}
{\large \ggn{29}.}
\end{center}\vspace*{-7pt}

{\greekfont Ἐὰν περισσὸς ἀριθμὸς περισσὸν ἀριθμὸν πολλαπλασιάς-ας ποιῇ
τινα, ὁ γενόμενος περισσὸς ἔσται.}

\epsfysize=0.8in
\centerline{\epsffile{Book09/fig28g.eps}}

{\greekfont Περισσὸς  γὰρ ἀριθμὸς ὁ Α περισσὸν τὸν Β πολλαπλασιάσας τὸν Γ
ποιείτω; λέγω, ὅτι ὁ Γ περισσός
ἐστιν.}

{\greekfont Ἐπεὶ γὰρ ὁ Α τὸν Β πολλαπλασιάσας τὸν Γ πεποίηκεν, ὁ Γ
ἄρα σύγκειται ἐκ τοσούτων ἴσων τῷ Β, ὅσαι εἰσὶν ἐν τῷ
Α μονάδες. καί ἐστιν ἑκάτερος τῶν Α, Β περισσός; ὁ Γ ἄρα σύγκειται
ἐκ  περισσῶν ἀριθμῶν, ὧν τὸ πλῆθος περισσόν ἐστιν. ὥστε
ὁ Γ περισσός ἐστιν; ὅπερ ἔδει δεῖχαι.}}

\ParallelRText{
\begin{center}
{\large Proposition 29}
\end{center}

If  an odd number makes some (number by)
multiplying an odd (number) then the created (number) will be odd.

\epsfysize=0.8in
\centerline{\epsffile{Book09/fig28e.eps}}

For let the odd number $A$ make $C$ (by) multiplying the odd (number)
$B$. I say that $C$ is odd.

For since $A$ has made $C$ (by) multiplying $B$, $C$ is thus composed  out of so many (magnitudes) equal to $B$, as many as (there) are  units in $A$ [Def.~7.15]. And each of $A$, $B$ is odd. Thus, $C$ is composed
out of odd (numbers), (and) the multitude of them is odd.   
Hence $C$ is odd [Prop.~9.23]. (Which is) the very thing it was required to show.}
\end{Parallel}

%%%%%%
% Prop 9.30
%%%%%%
\pdfbookmark[1]{Proposition 9.30}{pdf9.30}
\begin{Parallel}{}{} 
\ParallelLText{
\begin{center}
{\large \ggn{30}.}
\end{center}\vspace*{-7pt}

{\greekfont Ἐὰν περισσὸς ἀριθμὸς ἄρτιον ἀριθμὸν μετρῇ, καὶ τὸν ἥμισυν
α ὐτοῦ μετρήσει.}

\epsfysize=0.8in
\centerline{\epsffile{Book09/fig30g.eps}}

{\greekfont Περισσὸς γὰρ ἀριθμὸς ὁ Α ἄρτιον τὸν Β μετρείτω; λέγω, ὅτι
καὶ τὸν ἥμισυν α ὐτοῦ μετρήσει.}

{\greekfont Ἐπεὶ γὰρ ὁ Α τὸν Β μετρεῖ, μετρείτω α ὐτὸν κατὰ τὸν  Γ; λέγω,
ὅτι ὁ Γ οὐκ ἔστι περισσός. εἰ γὰρ δυνατόν, ἔστω. καὶ ἐπεὶ
ὁ Α τὸν Β μετρεῖ κατὰ τὸν Γ, ὁ Α ἄρα τὸν Γ πολλαπλασιάσας
τὸν Β πεποίηκεν. ὁ Β ἄρα σύγκειται ἐκ περισσῶν ἀριθμῶν,
ὧν τὸ πλῆθος περισσόν ἐστιν. ὁ Β ἄρα περισσός ἐστιν; ὅπερ
ἄτοπον; ὑπόκειται γὰρ ἄρτιος. οὐκ ἄρα ὁ Γ περισσός ἐστιν;
ἄρτιος ἄρα ἐστὶν ὁ Γ. ὥστε ὁ Α τὸν Β μετρεῖ ἀρτιάκις.
διὰ δὴ τοῦτο καὶ τὸν ἥμισυν α ὐτοῦ μετρήσει; ὅπερ ἔδει
δεῖχαι.}}

\ParallelRText{
\begin{center}
{\large Proposition 30}
\end{center}

If an odd number measures an even number then it
will also measure (one) half of it.

\epsfysize=0.8in
\centerline{\epsffile{Book09/fig30e.eps}}

For let the odd number $A$ measure the even (number) $B$. I say that
($A$) will also measure (one) half of ($B$).

For since $A$ measures $B$, let it measure it according to $C$. I say that
$C$ is not odd. For, if possible, let it be (odd). And since $A$ measures $B$ according to $C$, $A$ has thus made $B$ (by) multiplying $C$. Thus, $B$
is composed out of odd numbers, (and) the multitude of them is odd. $B$
is thus odd [Prop.~9.23]. The very thing (is)
absurd. For ($B$) was assumed (to be) even. Thus, $C$ is not odd.
Thus, $C$ is even. Hence, $A$ measures $B$ an even number of times.
So, on account of this, ($A$) will also measure (one) half of ($B$).
(Which is) the very thing it was required to show.}
\end{Parallel}

%%%%%%
% Prop 9.31
%%%%%%
\pdfbookmark[1]{Proposition 9.31}{pdf9.31}
\begin{Parallel}{}{} 
\ParallelLText{
\begin{center}
{\large \ggn{31}.}
\end{center}\vspace*{-7pt}

{\greekfont Ἐὰν περισσὸς ἀριθμὸς πρός τινα ἀριθμὸν πρῶτος ᾖ, καὶ πρὸς
τὸν διπλασίονα α ὐτοῦ πρῶτος ἔσται.}

\epsfysize=1.1in
\centerline{\epsffile{Book09/fig31g.eps}}

{\greekfont Περισσὸς γὰρ ἀριθμὸς ὁ Α πρός τινα ἀριθμὸν τὸν Β πρῶτος ἔστω,
τοῦ δὲ Β διπλασίων ἔστω ὁ Γ; λέγω, ὅτι ὁ Α [καὶ]
πρὸς τὸν Γ πρῶτός ἐστιν.}

{\greekfont Εἰ γὰρ μή εἰσιν [οἱ Α, Γ] πρῶτοι, μετρήσει τις α ὐτοὺς ἀριθμός.
μετρείτω, καὶ ἔστω ὁ Δ. καί ἐστιν ὁ Α περισσός; περισσὸς ἄρα καὶ
ὁ Δ. καὶ ἐπεὶ ὁ Δ περισσὸς ὢν τὸν Γ μετρεῖ, καί ἐστιν ὁ Γ
ἄρτιος, καὶ τὸν ἥμισυν ἄρα τοῦ  Γ μετρήσει [ὁ Δ]. τοῦ δὲ
Γ ἥμισύ ἐστιν ὁ Β; ὁ Δ ἄρα τὸν Β μετρεῖ. μετρεῖ δὲ καὶ τὸν Α.
ὁ Δ ἄρα τοὺς Α, Β μετρεῖ πρώτους ὄντας πρὸς ἀλλήλους;
ὅπερ ἐστὶν ἀδύνατον. οὐκ ἄρα ὁ Α πρὸς τὸν Γ πρῶτος οὔκ
ἐστιν. οἱ Α, Γ ἄρα πρῶτοι πρὸς ἀλλήλους εἰσίν; ὅπερ ἔδει
δεῖχαι.}}

\ParallelRText{
\begin{center}
{\large Proposition 31}
\end{center}

If an odd number is prime to some number then it
will also be prime to its double.

\epsfysize=1.1in
\centerline{\epsffile{Book09/fig31e.eps}}

For let the odd number $A$ be prime to some number $B$. And let $C$ be
double $B$. I say that $A$ is [also] prime to $C$.

For if [$A$ and $C$] are not prime (to one another) then some number
will measure them. Let it measure (them), and let it be $D$. And
$A$ is odd. Thus, $D$ (is) also odd. And since $D$, which is odd, measures
$C$, and $C$ is even, [$D$] will thus also measure half of $C$ [Prop.~9.30]. 
And $B$ is half of $C$.  Thus, $D$ measures $B$. And it also measures
$A$. Thus, $D$ measures (both) $A$ and $B$, (despite) them being prime to one another. The very thing is impossible. Thus, $A$ is not unprime 
to $C$. Thus, $A$ and $C$ are prime to one another. (Which is) the very thing it
was required to show.}
\end{Parallel}

%%%%%%
% Prop 9.32
%%%%%%
\pdfbookmark[1]{Proposition 9.32}{pdf9.32}
\begin{Parallel}{}{} 
\ParallelLText{
\begin{center}
{\large\ggn{32}.}
\end{center}\vspace*{-7pt}

{\greekfont Τῶν ἀπὸ δύαδος διπλασιαζομένων ἀριθμων ἕκαστος ἀρτιάκις
ἄρτιός ἐστι μόνον.}

\epsfysize=1.1in
\centerline{\epsffile{Book09/fig32g.eps}}

{\greekfont Ἀπὸ γὰρ δύαδος τῆς Α δεδιπλασιάσθωσαν ὁσοιδηποτοῦν ἀριθμοὶ
οἱ Β, Γ, Δ; λέγω, ὅτι οἱ Β, Γ, Δ ἀρτιάκις ἄρτιοί εἰσι μόνον.}

{\greekfont Ὅτι μὲν οὖν ἕκαστος [τῶν Β, Γ, Δ] ἀρτιάκις ἄρτιός
ἐστιν, φανερόν; ἀπὸ γὰρ δυάδος ἐστὶ διπλασιασθείς. λέγω, ὅτι
καὶ μόνον. ἐκκείσθω γὰρ μονάς. ἐπεὶ οὖν ἀπὸ μονάδος
ὁποσοιοῦν ἀριθμοὶ ἑχῆς ἀνάλογόν εἰσιν, ὁ δὲ μετὰ τὴν
μονάδα ὁ Α πρῶτός ἐστιν, ὁ μέγιστος τῶν Α, Β, Γ, Δ ὁ Δ ὑπ
οὐδενὸς ἄλλου μετρηθήσεται παρὲχ τῶν Α, Β, Γ. καί ἐστιν ἕκαστος
τῶν Α, Β, Γ ἄρτιος; ὁ Δ ἄρα ἀρτιάκις ἄρτιός ἐστι μόνον.
ὁμοίως δὴ δείχομεν, ὅτι [καὶ] ἑκάτερος τῶν Β, Γ
ἀρτιάκις ἄρτιός ἐστι μόνον; ὅπερ ἔδει δεῖχαι.}}

\ParallelRText{
\begin{center}
{\large Proposition 32}
\end{center}

Each of the numbers (which is continually) doubled,
(starting) from a dyad, is an even-times-even (number) only.

\epsfysize=1.1in
\centerline{\epsffile{Book09/fig32e.eps}}

For let any multitude of numbers whatsoever, $B$, $C$, $D$, be
(continually) doubled, (starting) from the dyad $A$. I say that
$B$, $C$, $D$ are  even-times-even (numbers) only.

In fact, (it is) clear that each [of $B$, $C$, $D$] is an even-times-even (number). For it is doubled from a dyad [Def.~7.8]. 
I also say that (they are even-times-even numbers) only. For let a unit
be laid down. Therefore, since any multitude of numbers whatsoever
are continuously proportional, starting from a unit, and the (number) $A$ after the unit is prime, the greatest of $A$, $B$, $C$, $D$, (namely) $D$,
will not be measured by any other (numbers) except $A$, $B$, $C$
[Prop.~9.13].
And each of $A$, $B$, $C$ is even. Thus, $D$ is an
even-time-even (number) only [Def.~7.8]. So, similarly,
we can show that each of $B$, $C$ is [also]  an even-time-even (number) only. (Which is) the very thing it was required to show.}
\end{Parallel}

%%%%%%
% Prop 9.33
%%%%%%
\pdfbookmark[1]{Proposition 9.33}{pdf9.33}
\begin{Parallel}{}{} 
\ParallelLText{
\begin{center}
{\large \ggn{33}.}
\end{center}\vspace*{-7pt}

{\greekfont Ἐὰν ἀριθμὸς τὸν ἥμισυν ἔθῃ περισσόν, ἀρτιάκις περισσός ἐστι
μόνον.}

\epsfysize=0.175in
\centerline{\epsffile{Book09/fig33g.eps}}

{\greekfont Ἀριθμὸς γὰρ ὁ Α τὸν ἥμισυν ἐθέτω περισσόν; λέγω, ὅτι ὁ Α
ἀρτιάκις περισσός ἐστι μόνον.}

{\greekfont Ὅτι μὲν οὖν ἀρτιάκις περισσός ἐστιν, φανερόν; ὁ γὰρ ἥμισυς
α ὐτοῦ περισσὸς ὢν μετρεῖ α ὐτὸν ἀρτιάκις, λέγω δή, ὅτι
καὶ μόνον. εἰ γὰρ ἔσται ὁ Α καὶ ἀρτιάκις ἄρτιος, μετρηθήσεται
ὑπὸ ἀρτίου κατὰ ἄρτιον ἀριθμόν; ὥστε καὶ ὁ ἥμισυς α ὐτοῦ
μετρηθήσεται ὑπὸ ἀρτίου ἀριθμοῦ περισσὸς ὤν; ὅπερ ἐστὶν
ἄτοπον. ὁ Α ἄρα ἀρτιάκις περισσός ἐστι μόνον; ὅπερ
ἔδει δεῖχαι.}}

\ParallelRText{
\begin{center}
{\large Proposition 33}
\end{center}

If a number has an odd half then it is an
even-time-odd (number) only.

\epsfysize=0.175in
\centerline{\epsffile{Book09/fig33e.eps}}

For let the number $A$ have an odd half. I say that $A$ is  an
even-times-odd (number) only.

In fact, (it is) clear that ($A$) is an even-times-odd (number). For its
half, being odd, measures it an even number of times [Def.~7.9]. So I also say that (it is an even-times-odd number) only. For if $A$ is also  an even-times-even
(number) then it will be measured by an even (number) according to
an even number [Def.~7.8]. Hence, its half
will also be measured by an even number, (despite)  being odd. The very
thing is absurd. Thus, $A$ is  an even-times-odd (number) only. (Which is)
the very thing it was required to show.}
\end{Parallel}

%%%%%%
% Prop 9.34
%%%%%%
\pdfbookmark[1]{Proposition 9.34}{pdf9.34}
\begin{Parallel}{}{} 
\ParallelLText{
\begin{center}
{\large \ggn{34}.}
\end{center}\vspace*{-7pt}

{\greekfont Ἐὰν ἀριθμὸς μήτε τῶν ἀπὸ δυάδος διπλασιαζομένων ᾖ, μήτε
τὸν ἥμισυν ἔθῃ περισσόν, ἀρτιάκις τε ἄρτιός ἐστι καὶ ἀρτιάκις
περισσός.}

\epsfysize=0.175in
\centerline{\epsffile{Book09/fig33g.eps}}

{\greekfont Ἀριθμὸς γὰρ ὁ Α μήτε τῶν ἀπὸ δυάδος διπλασιαζομένων ἔστω
μήτε τὸν ἥμισυν ἐθέτω περισσόν; λέγω, ὅτι ὁ Α ἀρτιάκις
τέ ἐστιν ἄρτιος καὶ ἀρτιάκις περισσός.}

{\greekfont Ὅτι μὲν οὖν ὁ Α ἀρτιάκις ἐστὶν ἄρτιος, φανερόν; τὸν γὰρ
ἥμισυν οὐκ ἔχει περισσόν. λέγω δή, ὅτι καὶ ἀρτιάκις περισσός
ἐστιν. ἐὰν γὰρ τὸν Α τέμνωμεν δίθα καὶ τὸν ἥμισυν α ὐτοῦ
δίθα καὶ τοῦτο ἀεὶ ποιῶμεν, καταντήσομεν εἴς τινα
ἀριθμὸν περισσόν, ὃς μετρήσει τὸν Α κατὰ ἄρτιον ἀριθμόν. εἰ γὰρ
οὔ, καταντήσομεν εἰς
δυάδα, καὶ
ἔσται ὁ Α τῶν ἀπὸ δυάδος διπλασιαζομένων; ὅπερ οὐθ
ὑπόκειται. ὥστε ὁ Α ἀρτιάκις περισσόν ἐστιν. ἐδείθθη δὲ καὶ
ἀρτιάκις ἄρτιος. ὁ Α ἄρα ἀρτιάκις τε ἄρτιός ἐστι καὶ ἀρτιάκις
περισσός; ὅπερ ἔδει δεῖχαι.}}

\ParallelRText{
\begin{center}
{\large Proposition 34}
\end{center}

If a number is neither (one) of the (numbers) doubled from a dyad, nor
has an odd half, then it is (both) an even-times-even and an even-times-odd
(number).

\epsfysize=0.175in
\centerline{\epsffile{Book09/fig33e.eps}}

For let the number $A$   neither be (one) of the (numbers) doubled from a
dyad, nor let it have an odd half. I say that $A$ is (both) an even-times-even
and an even-times-odd (number).

In fact, (it is) clear that $A$ is an even-times-even (number) [Def.~7.8]. For it does
not have an odd half. So I say that it is also an even-times-odd (number). 
For if we cut $A$ in half, and (then cut) its half in half, and we do this continually, then we will arrive at some odd number which will measure
$A$ according to an even number. For if not, we will arrive at a dyad,
and $A$ will be (one) of the (numbers) doubled from a dyad. The very opposite thing (was) assumed. Hence, $A$ is an even-times-odd (number)
[Def.~7.9]. And it was also shown (to be) an
even-times-even (number). Thus, $A$ is (both) an even-times-even and
an even-times-odd (number). (Which is) the very thing it was required to
show.}
\end{Parallel}

%%%%%%
% Prop 9.35
%%%%%%
\pdfbookmark[1]{Proposition 9.35}{pdf9.35}
\begin{Parallel}{}{} 
\ParallelLText{
\begin{center}
{\large \ggn{35}.}
\end{center}\vspace*{-7pt}

{\greekfont Ἐὰν ὦσιν ὁσοιδηποτοῦν ἀριθμοὶ ἑχῆς ἀνάλογον, ἀφαιρεθῶσι δὲ ἀπό τε τοῦ δευτέρου καὶ τοῦ ἐσθάτου ἴσοι τῷ πρώτῳ, ἔσται
ὡς ἡ τοῦ δευτέρου ὑπεροθὴ πρὸς τὸν πρῶτον, οὕτως ἡ τοῦ
ἐσθάτου ὑπεροθὴ πρὸς τοὺς πρὸ ἑαυτοῦ πάντας.}\\

\epsfysize=1.6in
\centerline{\epsffile{Book09/fig35g.eps}}

{\greekfont Ἔστωσαν ὁποσοιδηποτοῦν ἀριθμοὶ ἑχῆς ἀνάλογον
οἱ Α, ΒΓ, Δ, ΕΖ ἀφθόμενοι ἀπὸ ἐλαθίστου τοῦ Α, καὶ ἀφῃρήσθω ἀπὸ τοῦ ΒΓ καὶ τοῦ ΕΖ τῲ Α ἴσος ἑκάτερος τῶν ΒΗ, ΖΘ; λέγω, ὅτι ἐστὶν ὡς ὁ ΗΓ πρὸς τὸν Α, οὕτως ὁ ΕΘ πρὸς τοὺς Α, ΒΓ, Δ.}

{\greekfont Κείσθω γὰρ τῷ μὲν ΒΓ ἴσος ὁ ΖΚ, τῷ δὲ Δ ἴσος ὁ  ΖΛ. καὶ
ἐπεὶ ὁ ΖΚ τῷ ΒΓ ἴσος
 ἐστίν, ὧν ὁ ΖΘ τῷ ΒΗ ἴσος ἐστίν, λοιπὸς ἄρα ὁ ΘΚ λοιπῷ
 τῷ ΗΓ ἐστιν ἴσος.
καὶ ἐπεί ἐστιν ὡς ὁ ΕΖ πρὸς τὸν Δ, οὕτως ὁ Δ πρὸς τὸν ΒΓ καὶ
ὁ ΒΓ πρὸς τὸν Α, ἴσος δὲ ὁ μὲν Δ τῷ ΖΛ, ὁ δὲ ΒΓ τῷ ΖΚ, ὁ
δὲ Α τῷ ΖΘ, ἔστιν ἄρα ὡς ὁ ΕΖ πρὸς τὸν ΖΛ, οὕτως ὁ ΛΖ πρὸς
τὸν ΖΚ καὶ ὁ ΖΚ πρὸς τὸν ΖΘ. διελόντι, ὡς ὁ ΕΛ πρὸς τὸν ΛΖ, οὕτως ὁ ΛΚ πρὸς τὸν ΖΚ καὶ ὁ ΚΘ πρὸς τὸν ΖΘ. ἔστιν ἄρα καὶ
ὡς εἷς τῶν ἡγουμένων πρὸς ἕνα τῶν ἑπομένων, οὕτως
ἅπαντες οἱ ἡγούμενοι πρὸς ἅπαντας τοὺς ἑπομένους; ἔστιν
ἄρα ὡς ὁ ΚΘ πρὸς τὸν ΖΘ, οὕτως οἱ ΕΛ, ΛΚ, ΚΘ πρὸς τοὺς ΛΖ, ΖΚ,
ΘΖ.  ἴσος δὲ ὁ μὲν ΚΘ τῷ ΓΗ, ὁ δὲ ΖΘ τῷ Α, οἱ δὲ ΛΖ, ΖΚ, ΘΖ
τοὶς Δ, ΒΓ, Α; ἔστιν ἄρα ὡς ὁ ΓΗ  πρὸς τὸν  Α, οὕτως ὁ ΕΘ πρὸς τοὺς Δ, ΒΓ, Α. ἔστιν ἄρα ὡς ἡ τοῦ δευτέρου ὑπεροθὴ πρὸς
τὸν πρῶτον, οὕτως ἡ τοῦ
ἐσθάτου ὑπεροθὴ πρὸς τοὺς πρὸ ἑαυτοῦ πάντας;
ὅπερ ἔδει δεῖχαι.}}

\ParallelRText{
\begin{center}
{\large Proposition 35}$^\dag$
\end{center}

If there is any multitude whatsoever of continually
proportional numbers, and (numbers) equal to the first are subtracted from (both) the
second and the last, then as the excess of the second (number is) to the first, so
the excess of the last will be to (the sum of) all   those (numbers) before it.

\epsfysize=1.6in
\centerline{\epsffile{Book09/fig35e.eps}}

Let $A$, $BC$, $D$, $EF$ be any multitude whatsoever of continuously proportional numbers, beginning from the least $A$.  And let $BG$ and
$FH$, each equal to $A$, be subtracted from $BC$ and $EF$ (respectively). I say that as $GC$ is to $A$, so $EH$ is to $A$, $BC$, $D$.

For let $FK$ be made equal to $BC$, and $FL$ to  $D$. And since $FK$
is equal to $BC$, of which $FH$ is equal to $BG$, the remainder,
$HK$,  is thus equal to the remainder, $GC$. And since as $EF$ is to $D$, so
$D$ (is) to $BC$, and $BC$ to $A$ [Prop.~7.13],
and $D$ (is) equal to $FL$, and $BC$ to $FK$, and $A$ to  $FH$, thus
as $EF$ is to $FL$, so $LF$ (is) to $FK$, and $FK$ to $FH$. By
separation, as $EL$ (is) to $LF$, so $LK$ (is) to $FK$, and $KH$ to $FH$
[Props.~7.11, 7.13]. And thus
as one of the leading (numbers) is to one of the following, so (the sum of) all of the leading (numbers is) to (the sum of) all of the following [Prop.~7.12]. Thus, as $KH$ is to $FH$, so
$EL$, $LK$, $KH$ (are) to $LF$, $FK$, $HF$.  And $KH$ (is) equal to
$CG$, and $FH$ to $A$, and $LF$, $FK$, $HF$ to $D$, $BC$, $A$.
Thus, as $CG$ is to $A$, so $EH$ (is) to $D$, $BC$, $A$. Thus,
as the excess of the second (number) is to the first, so the excess of the last
(is) to (the sum of) all  those (numbers) before it. (Which is) the very thing it was required to show.}
\end{Parallel}
{\footnotesize\noindent$^\dag$ This proposition allows us to sum a geometric series of the form $a$, $a\,r$, $a\,r^2$, $a\,r^3,\cdots a\,r^{n-1}$. 
According to Euclid, the sum, $S_n$,  satisfies
$(a\,r-a)/a = (a\,r^n-a)/S_n$. Hence, $S_n= a\,(r^n-1)/(r-1)$.}

%%%%%%
% Prop 9.36
%%%%%%
\pdfbookmark[1]{Proposition 9.36}{pdf9.36}
\begin{Parallel}{}{} 
\ParallelLText{
\begin{center}
{\large \ggn{36}.}
\end{center}\vspace*{-7pt}

{\greekfont Ἐὰν ἀπὸ μονάδος ὁποσοιοῦν ἀριθμοὶ ἑχῆς ἐκτεθῶσιν ἐν τῇ
διπλασίονι ἀναλογίᾳ, ἕως οὗ ὁ σύμπας συντεθεὶς πρῶτος
γένηται, καὶ ὁ σύμπας ἐπὶ τὸν ἔσθατον πολλαπλασιασθεὶς
ποιῇ τινα, ὁ γενόμενος τέλειος ἔσται.}

{\greekfont Ἀπὸ γὰρ μονάδος ἐκκείσθωσαν ὁσοιδηποτοῦν ἀριθμ-οὶ
ἐν τῇ διπλασίονι ἀναλογίᾳ, ἕως οὗ ὁ σύμπας συντεθεὶς
πρῶτος γένηται, οἱ Α, Β, Γ, Δ, καὶ τῷ σύμπαντι ἴσος
ἔστω ὁ Ε, καὶ ὁ Ε τὸν Δ πολλαπλασιάσας τὸν ΖΗ ποιείτω. λέγω,
ὅτι ὁ ΖΗ τέλειός ἐστιν.}\\

\epsfysize=1.2in
\centerline{\epsffile{Book09/fig36g.eps}}}

\ParallelRText{
\begin{center}
{\large Proposition 36}$^\dag$
\end{center}

If any multitude whatsoever of numbers is set out
continuously 
in a double  proportion, (starting) from a  unit,
until the whole sum added together becomes  prime, and the sum multiplied
into the last (number) makes some (number), then the (number so) created
will be perfect.

For let any multitude of numbers, $A$, $B$, $C$, $D$, be set out (continuouly) in a
double proportion, until the whole sum added together is made  prime. And let $E$ be equal to the sum. And let $E$ make $FG$ (by) multiplying $D$. I say that $FG$ is a perfect (number).

\epsfysize=1.3in
\centerline{\epsffile{Book09/fig36e.eps}}
}
\end{Parallel}~\\

\begin{Parallel}{}{} 
\ParallelLText{
\epsfysize=1.4in
\centerline{\epsffile{Book09/fig36ag.eps}}

{\greekfont Ὅσοι γάρ εἰσιν οἱ Α, Β, Γ, Δ τῷ πλήθει, τοσοῦτοι
ἀπὸ τοῦ Ε εἰλήφθωσαν ἐν τῇ διπλασίονι ἀναλογίᾳ οἱ 
Ε, ΘΚ, Λ, Μ; δι ἴσου ἄρα ἐστὶν ὡς ὁ Α πρὸς τὸν Δ, οὕτως
ὁ Ε πρὸς τὸν Μ. ὁ ἄρα ἐκ τῶν Ε, Δ ἴσος ἐστὶ τῷ ἐκ τῶν
Α, Μ. καί ἐστιν ὁ ἐκ τῶν Ε, Δ 
ὁ ΖΗ; καὶ ὁ ἐκ τῶν
Α, Μ ἄρα ἐστὶν ὁ ΖΗ. ὁ Α ἄρα τὸν Μ πολλαπλασιάσας τὸν ΖΗ
πεποίηκεν; ὁ Μ ἄρα τὸν ΖΗ μετρεῖ κατὰ τὰς ἐν τῷ Α μονάδας.
καί ἐστι δυὰς ὁ Α; διπλάσιος  ἄρα ἐστὶν ὁ ΖΗ τοῦ Μ.
εἰσὶ δὲ καὶ οἱ Μ, Λ, ΘΚ, Ε ἑχῆς διπλάσιοι ἀλλήλων; οἱ
Ε, ΘΚ, Λ, Μ, ΖΗ ἄρα ἑχῆς ἀνάλογόν εἰσιν ἐν τῇ διπλασίονι ἀναλογίᾳ. ἀφῃρήσθω δὴ ἀπὸ τοῦ δευτέρου τοῦ ΘΚ καὶ τοῦ
ἐσθάτου τοῦ ΖΗ τῷ πρώτῳ τῷ Ε ἴσος ἑκάτερος τῶν ΘΝ, ΖΧ;
ἔστιν ἄρα ὡς ἡ τοῦ δευτέρου ἀριθμοῦ ὑπεροθὴ πρὸς
τὸν πρῶτον, οὕτως ἡ τοῦ ἐσθάτου ὑπεροθὴ πρὸς τοὺς πρὸ
ἑαυτοῦ πάντας. ἔστιν ἄρα ὡς ὁ ΝΚ πρὸς τὸν Ε, οὕτως
ὁ ΧΗ πρὸς τοὺς Μ, Λ, ΚΘ, Ε. καί ἐστιν ὁ ΝΚ ἴσος τῷ Ε;
καὶ ὁ ΧΗ ἄρα ἴσος ἐστὶ τοῖς Μ, Λ, ΘΚ, Ε. ἔστι δὲ καὶ
ὁ ΖΧ τῷ Ε ἴσος, ὁ δὲ Ε τοῖς Α, Β, Γ, Δ καὶ τῇ μονάδι. ὅλος
ἄρα ὁ ΖΗ ἴσος ἐστὶ τοῖς τε Ε, ΘΚ, Λ, Μ καὶ τοῖς Α, Β, Γ, Δ καὶ τῇ
μονάδι; καὶ μετρεῖται ὑπ α ὐτῶν. λέγω, ὅτι καὶ ὁ ΖΗ ὐπ οὐδενὸς ἄλλου μετρηθήσεται παρὲχ τῶν Α, Β, Γ, Δ, Ε, ΘΚ, Λ, Μ καὶ
τῆς μονάδος. εἰ γὰρ δυνατόν, μετρείτω τις τὸν ΖΗ ὁ Ο, καὶ ὁ Ο
μ ηδενὶ τῶν Α, Β, Γ, Δ, Ε, ΘΚ, Λ, Μ ἔστω ὁ α ὐτός. καὶ ὁσάκις
ὁ Ο τὸν ΖΗ μετρεῖ, τοσαῦται μονάδες ἔστωσαν ἐν τῷ Π; ὁ Π
ἄρα τὸν Ο πολλαπλασιάσας τὸν ΖΗ πεποίηκεν. ἀλλὰ μὴν καὶ ὁ Ε
τὸν Δ πολλαπλασιάσας τὸν ΖΗ πεποίηκεν; ἔστιν ἄρα ὡς ὁ Ε πρὸς
τὸν Π, ὁ Ο πρὸς τὸν Δ. καὶ ἐπεὶ ἀπὸ μονάδος ἑχῆς
ἀνάλογόν εἰσιν οἱ Α, Β, Γ, Δ, ὁ Δ ἄρα ὑπ οὐδενὸς ἄλλου
ἀριθμοῦ μετρηθήσεται παρὲχ τῶν Α, Β, Γ. καὶ ὑπόκειται ὁ Ο οὐδενὶ
τῶν Α, Β, Γ ὁ α ὐτός; οὐκ ἄρα μετρήσει ὁ Ο τὸν Δ. ἀλλ ὡς ὁ Ο
πρὸς τὸν Δ, ὁ Ε πρὸς τὸν Π; οὐδὲ ὁ Ε ἄρα τὸν Π μετρεῖ. καί
ἐστιν ὁ Ε πρῶτος; πᾶς δὲ πρῶτος ἀριθμὸς πρὸς
ἅπαντα, ὃν μὴ μετρεῖ, πρῶτός [ἐστιν]. οἱ Ε, Π
ἄρα πρῶτοι πρὸς ἀλλήλους εἰσίν. οἱ δὲ πρῶτοι καὶ
ἐλάθιστοι, οἱ δὲ ἐλάθιστοι μετροῦσι τοὺς τὸν α ὐτὸν λόγον ἔθοντας
ἰσάκις  ὅ τε ἡγούμενος τὸν  ἡγούμενον καὶ ὁ  ἑπόμενος τὸν
ἑπόμενον; καί ἐστιν ὡς ὁ Ε πρὸς τὸν Π, ὁ Ο
πρὸς τὸν Δ;  ἰσάκις ἄρα ὁ Ε τὸν Ο μετρεῖ καὶ ὁ Π τὸν Δ.
 ὁ δὲ Δ ὑπ οὐδενὸς ἄλλου μετρεῖται
παρὲχ τῶν Α, Β, Γ; ὁ Π ἄρα ἑνὶ τῶν Α, Β, Γ ἐστιν ὁ α ὐτός. ἔστω
τῷ Β ὁ α ὐτός. καὶ ὅσοι εἰσὶν οἱ Β, Γ, Δ τῷ πλήθει τοσοῦτοι
εἰλήφθωσαν ἀπὸ τοῦ Ε οἱ Ε, ΘΚ, Λ. καί εἰσιν οἱ Ε, ΘΚ, Λ τοῖς
Β, Γ, Δ ἐν τῷ α ὐτῷ λόγω; δι ἴσου ἄρα ἐστὶν ὡς ὁ Β πρὸς
τὸν Δ, ὁ Ε πρὸς τὸν Λ. ὁ ἄρα ἐκ τῶν Β, Λ ἴσος ἐστὶ τῷ ἐκ τῶν
Δ, Ε; ἀλλ ὁ ἐκ τῶν Δ, Ε ἴσος ἐστὶ τῷ ἐκ τῶν
Π, Ο; καὶ ὁ ἐκ τῶν Π, Ο ἄρα ἴσος ἐστὶ τῷ ἐκ τῶν Β, Λ. ἔστιν
ἄρα ὡς ὁ Π πρὸς τὸν Β, ὁ Λ πρὸς τὸν Ο. καί ἐστιν ὁ Π τῷ Β
ὁ α ὐτός; καὶ ὁ Λ ἄρα τῳ Ο ἐστιν ὁ α ὐτός;  ὅπερ ἀδύνατον;
ὁ γὰρ Ο ὑπόκειται μ ηδενὶ τῶν ἐκκειμένων ὁ α ὐτός; οὐκ
ἄρα τὸν ΖΗ μετρήσει τις ἀριθμὸς παρὲχ τῶν Α, Β, Γ, Δ, Ε, ΘΚ, Λ, Μ
καὶ τῆς μονάδος. καὶ ἐδείθη ὁ ΖΗ τοῖς Α, Β, Γ, Δ, Ε, ΘΚ, Λ, Μ
καὶ τῇ μονάδι ἴσος. τέλειος δὲ ἀριθμός ἐστιν ὁ τοῖς ἑαυτοῦ
μέρεσιν ἴσος ὤν; τέλειος ἄρα ἐστὶν ὁ ΖΗ; ὅπερ ἔδει δεῖχαι.}}

\ParallelRText{
\epsfysize=1.4in
\centerline{\epsffile{Book09/fig36ae.eps}}

For as many as is the multitude of $A$, $B$, $C$, $D$, let so many (numbers), $E$, $HK$, $L$, $M$, be taken in a double
proportion, (starting) from $E$. Thus, via equality, as $A$ is to $D$, so
$E$ (is) to $M$ [Prop.~7.14].
Thus, the (number created) from (multiplying) $E, D$ is equal to the (number
created) from (multiplying) $A$, $M$.
 And $FG$ is the
(number created) from (multiplying) $E$, $D$. Thus, $FG$ is also
the (number created) from (multiplying) $A$, $M$ [Prop.~7.19]. Thus, $A$ has made $FG$ (by) multiplying $M$. Thus, $M$ measures $FG$ according to the units in $A$.
And $A$ is a dyad. Thus, $FG$ is double $M$.  And $M$, $L$, $HK$,
$E$ are also continuously double one another. Thus, $E$, $HK$, $L$, $M$,
$FG$ are continuously proportional in a double proportion. So
let $HN$ and $FO$, each equal to the first (number) $E$, be subtracted from the second (number) $HK$ and the last $FG$ (respectively). 
Thus, as the excess of the second number is to the first, so the excess
of the last (is) to (the sum of) all  those (numbers) before it [Prop.~9.35]. Thus, as $NK$ is to $E$, so
$OG$ (is) to $M$, $L$, $KH$, $E$. And $NK$ is equal to $E$. And thus
$OG$ is equal to $M$, $L$, $HK$, $E$. And $FO$ is also equal to $E$,
and $E$ to $A$, $B$, $C$, $D$, and a unit. Thus, the whole of $FG$ is equal to
$E$, $HK$, $L$, $M$, and $A$, $B$, $C$, $D$, and a unit. And it is
measured by them. I also say that $FG$ will be measured by no other (numbers)
except $A$, $B$, $C$, $D$, $E$, $HK$, $L$, $M$, and a unit. For, if possible, let some (number) $P$ measure $FG$, and let $P$ not be the same
as any of $A$, $B$, $C$, $D$, $E$, $HK$, $L$, $M$. And as many times
as $P$ measures $FG$, so many units let there be in $Q$. Thus, $Q$ has
made $FG$ (by) multiplying $P$. But, in fact, $E$ has  also made $FG$
(by) multiplying $D$. Thus, as $E$ is to $Q$, so $P$ (is) to $D$
[Prop.~7.19]. And since $A$, $B$, $C$, $D$
are continually proportional, (starting) from a unit, $D$ will thus not be measured
by any other numbers except $A$, $B$, $C$ [Prop.~9.13]. And $P$ was assumed not (to be) the same as any of $A$, $B$, $C$. Thus, $P$ does not measure $D$. But, as $P$ (is) to
$D$, so $E$ (is) to $Q$. Thus, $E$ does not measure $Q$ either
 [Def.~7.20]. And $E$ is a prime (number).
And every prime number [is] prime to every  (number) which it does not
measure [Prop.~7.29]. Thus, $E$ and $Q$ are
prime to one another. And  (numbers) prime (to one another are) also the least (of those
numbers having the same ratio as them) [Prop.~7.21],
and the least (numbers) measure those (numbers) having the same ratio as them an equal number of times, the leading (measuring) the leading, and the
following the following [Prop.~7.20].
And as $E$ is to $Q$, (so) $P$ (is) to $D$. Thus, $E$ measures $P$ the same number of times as $Q$ (measures) $D$. And $D$ is not measured by any
other (numbers) except $A$, $B$, $C$.  Thus, $Q$ is the same as one
of $A$, $B$, $C$. Let it be the same as $B$. And as many as is the
multitude of $B$, $C$, $D$, let so many (of the set out numbers) be taken, (starting) from $E$, 
(namely) $E$, $HK$, $L$. And $E$, $HK$, $L$ are in the same ratio as
$B$, $C$, $D$. Thus, via equality, as $B$ (is) to $D$, (so) $E$ (is) to $L$
[Prop.~7.14]. Thus, the (number created) from
(multiplying) $B$, $L$ is equal to the (number created) from multiplying
$D$, $E$ [Prop.~7.19]. But, the (number created)
from (multiplying) $D$, $E$ is equal to the (number created) from (multiplying)
$Q$, $P$. Thus, the (number created) from (multiplying) $Q$, $P$
is equal to the (number created) from (multiplying) $B$, $L$. Thus, as
$Q$ is to $B$, (so) $L$ (is) to $P$ [Prop.~7.19]. 
And $Q$ is the same as $B$. Thus, $L$ is also the same as $P$. The
very thing (is) impossible. For $P$ was assumed not (to be) the same
as any of the (numbers) set out. Thus, $FG$ cannot be measured by any number except $A$, $B$, $C$, $D$, $E$, $HK$, $L$, $M$, and a unit.
And $FG$ was shown (to be) equal to (the sum of) $A$, $B$, $C$, $D$, $E$, $HK$,
$L$, $M$, and a unit. And a perfect number is one which is equal to (the sum of) its own
parts [Def.~7.22]. Thus, $FG$ is a perfect (number).
(Which is) the very thing it was required to show.}
\end{Parallel}
{\footnotesize\noindent$^\dag$ This proposition demonstrates that perfect
numbers take the form $2^{n-1}\,(2^n-1)$ provided that $2^n-1$ is a prime
number. The ancient Greeks knew of four perfect numbers: 6, 28, 496, and
8128, which correspond to $n= 2$, 3, 5, and 7, respectively.}
